\section{Meromorphic functions and pseudo-morphisms}
\egalabel{IV.20}{20}

\setcounter{subsection}{-1}
\subsection{Introduction}
\EGAReaderTarget{ega4.c04.subsec.20-0}{IV.20.0}

Most of the notions and results of \textsection\textsection20 and 21 are directly related to Chapter~I, and hardly depend on Chapters~II to IV, except as regards the occasional use of the notions of depth and of regular local ring (in (20.6), \sref{IV.21.11}, \sref{IV.21.13} and \sref{IV.21.15}), of Zariski's ``main theorem'' in (20.4) and \sref{IV.21.12}, and of the properties of transversally regular immersions in (20.6) and \sref{IV.21.15}.

\oldpage[IV-4]{226}
In \textsection20, we introduce several variants of the notion of rational map, already studied in \sref[I]{section:I.7} from a point of view still rather close to the classical point of view, and for this reason rather ill suited to the case of preschemes that are not necessarily reduced.
The notions and results of \textsection20 are used to develop in \textsection21 (n\textsuperscript{os} \sref{IV.21.1} to \sref{IV.21.7}) the general notion of divisor and its most elementary properties.
The notion is particularly convenient when the local rings of the preschemes under consideration are Noetherian and integrally closed, and above all when they are moreover \emph{factorial} (\sref{IV.21.6} and \sref{IV.21.7}), because of its identification, in the latter case, with the notion of \emph{$1$-codimensional cycle} (linear combination of irreducible sub-preschemes of codimension~$1$).
In \sref{IV.21.9}, we determine the divisors on a Noetherian prescheme of dimension~$1$ that is not necessarily normal, which is useful for various applications.
N\textsuperscript{os} \sref{IV.21.11} and \sref{IV.21.12} give two important theorems, due respectively to Auslander--Buchsbaum and to Van der Waerden, and related to the notion of factorial ring (n\textsuperscript{os} \sref{IV.21.9}, \sref{IV.21.11} and \sref{IV.21.12} are independent of one another).
In n\textsuperscript{os} \sref{IV.21.13} and \sref{IV.21.14}, which are likewise independent of the three preceding ones, we study a useful variant of the notion of factorial local ring, namely that of \emph{parafactorial} local ring, which arises in particular~\cite{IV-41} in the development of comparison theorems between the Picard group of a projective prescheme $X$ over a field $k$ and that of a ``hyperplane section''.
We will see in \sref{IV.21.14.1} (Ramanujam--Samuel theorem) that parafactorial local rings are far more numerous than one could have expected \emph{a priori}.

In (20.5), (20.6) and \sref{IV.21.15}, we take up the preceding notions again, but from a point of view ``relative'' to a fixed base prescheme.
For the moment these notions are used only rather rarely; in particular, the notion of relative divisor is hardly used except for positive divisors, and in that case it is advantageously made explicit without the help of the notion of relative meromorphic function, by means of the notion of transversally regular immersion of codimension~$1$.
It will therefore be advantageous to omit these sections on a first reading.

\subsection{Meromorphic functions}
\EGAReaderTarget{ega4.c04.subsec.20-1}{IV.20.1}

\begin{env}[20.1.1]\label{IV.20.1.1}
\EGAReaderTarget{ega4.c04.statement.20-1-1}{IV.20.1.1}
Let $(X,\sh{O}_X)$ be a ringed space, and let $\sh{S}$ be a sub-sheaf \emph{of sets} of $\sh{O}_X$.
For every open subset $U$ of $X$, consider the \emph{ring of fractions} $\Gamma(U,\sh{O}_X)[\Gamma(U,\sh{S})^{-1}]$ (Bourbaki, \emph{Alg. comm.}, chap.~II, \textsection2, n\textsuperscript{o}1).
It is immediate that the map $U\mapsto\Gamma(U,\sh{O}_X)[\Gamma(U,\sh{S})^{-1}]$ is a \emph{presheaf of rings} (\sref[0]{0.1.5.1} and \sref[0]{0.1.5.7}).
We denote by $\sh{O}_X[\sh{S}^{-1}]$ the \emph{sheaf of rings} associated to this presheaf, and we say that it is the \emph{sheaf of rings of fractions of $\sh{O}_X$ with denominators in $\sh{S}$}; it is a \emph{flat} $\sh{O}_X$-module.
It is immediate that for every $x\in X$ we have a canonical isomorphism
\[
\label{IV.20.1.1.1}
  (\sh{O}_X[\sh{S}^{-1}])_x\isoto\sh{O}_{X,x}[\sh{S}_x^{-1}],
  \tag{20.1.1.1}
\]
for the reasoning of \sref[0]{0.1.4.5} generalizes immediately to the case where we have an inductive system $(A_\alpha,\vphi_{\beta\alpha})$ of rings and, for each index $\alpha$, a subset $S_\alpha$ of $A_\alpha$ such that
\oldpage[IV-4]{227}
$\vphi_{\beta\alpha}(S_\alpha)\subset S_\beta$ for $\alpha\leq\beta$; we then take for $S$ the inductive limit in $A=\varinjlim A_\alpha$ of the inductive system of subsets $(S_\alpha)$.
\end{env}

\begin{env}[20.1.2]\label{IV.20.1.2}
\EGAReaderTarget{ega4.c04.statement.20-1-2}{IV.20.1.2}
Now let $\sh{F}$ be an $\sh{O}_X$-module.
We then set
\[
\label{IV.20.1.2.1}
  \sh{F}[\sh{S}^{-1}]=\sh{F}\otimes_{\sh{O}_X}\sh{O}_X[\sh{S}^{-1}]
  \tag{20.1.2.1}
\]
and we say that it is the \emph{sheaf of modules of fractions of $\sh{F}$ with denominators in $\sh{S}$}; it is immediate that it is associated to the presheaf of modules $U\mapsto\Gamma(U,\sh{F})[\Gamma(U,\sh{S})^{-1}]$, and that for every $x\in X$ we have a canonical isomorphism
\[
\label{IV.20.1.2.2}
  (\sh{F}[\sh{S}^{-1}])_x\isoto\sh{F}_x[\sh{S}_x^{-1}].
  \tag{20.1.2.2}
\]
\end{env}

\begin{env}[20.1.3]\label{IV.20.1.3}
\EGAReaderTarget{ega4.c04.statement.20-1-3}{IV.20.1.3}
We will be interested here in the case where $\sh{S}$ is the sub-sheaf $\sh{S}(\sh{O}_X)$ of $\sh{O}_X$ such that, for every open subset $U$, $\Gamma(U,\sh{S})$ is the \emph{set of regular elements} of the ring $\Gamma(U,\sh{O}_X)$; it is immediate that this is a sheaf (and not only a presheaf), the regularity of a section of $\sh{O}_X$ over $U$ being checked ``stalk by stalk'' (i.e. meaning that the germ of the section at $x$ is regular in $\sh{O}_{X,x}$ for every $x\in U$); in other words, $\sh{S}(\sh{O}_X)_x$ is nothing other than the set of regular elements of $\sh{O}_{X,x}$.
The corresponding sheaf of rings
\[
  \sh{M}_X=\sh{O}_X[\sh{S}^{-1}]
\]
is called the \emph{sheaf of germs of meromorphic functions on $X$}, and the sections of $\sh{M}_X$ over $X$ are called the \emph{meromorphic functions on $X$}; they form a ring, which we denote by $M(X)$.
For every $\sh{O}_X$-module $\sh{F}$,
\[
  \sh{F}\otimes_{\sh{O}_X}\sh{M}_X=\sh{F}[\sh{S}^{-1}]
\]
is also denoted by $\sh{M}_X(\sh{F})$ and called the \emph{sheaf of germs of meromorphic sections of $\sh{F}$}; its sections over $X$ form an $M(X)$-module denoted by $M(X,\sh{F})$, whose elements are called \emph{meromorphic sections of $\sh{F}$ over $X$}.
These definitions imply that for every open subset $U$ of $X$ we have a canonical isomorphism $\sh{M}_X(\sh{F})|U\isoto\sh{M}_U(\sh{F}|U)$, in particular $\sh{M}_X|U\isoto\sh{M}_U$.
\end{env}

\begin{env}[20.1.3.1]
\label{IV.20.1.3.1}
If $X$ is a \emph{reduced prescheme}, note that if an element $s\in\Gamma(U,\sh{O}_X)$ is such that $s_\xi\neq0$ for every maximal point $\xi$ of $U$, then $s$ is \emph{regular}.
Indeed, if $st=0$ for some $t\in\Gamma(U,\sh{O}_X)$, we have $s_\xi t_\xi=0$, hence $t_\xi=0$ since $\sh{O}_{X,\xi}$ is a field, and to say that $t_\xi=0$ for every maximal point $\xi$ of $X$ means that $t=0$: indeed, we are immediately reduced to the case where $U$ is affine, and an element of a reduced ring that belongs to all the minimal prime ideals is zero by definition.
The converse is true if the set of irreducible components of $X$ is \emph{locally finite}.
Indeed, we are immediately reduced to the case where $X=\Spec(A)$ is affine; if $\mathfrak p_i$ ($1\leq i\leq n$) are the minimal prime ideals of $A$ and if $s\in\mathfrak p_i$ for some index $i$, there exists $t\in A$ such that $t\in\mathfrak p_j$ for $j\neq i$ and $t\notin\mathfrak p_i$ (Bourbaki, \emph{Alg. comm.}, chap.~II, \textsection1, n\textsuperscript{o}1, prop.~1); we therefore have $st\in\mathfrak p_i$ for every $i$, and consequently $st=0$ since $A$ is reduced; $s$ is therefore not regular.
\end{env}

\begin{env}[20.1.4]\label{IV.20.1.4}
\EGAReaderTarget{ega4.c04.statement.20-1-4}{IV.20.1.4}
For every open subset $U$ of $X$, the homomorphism $t\mapsto t/1$ from $\Gamma(U,\sh{O}_X)$ to $\Gamma(U,\sh{O}_X)[\Gamma(U,\sh{S})^{-1}]$ (which is nothing other than the \emph{total ring of fractions} of
\oldpage[IV-4]{228}
$\Gamma(U,\sh{O}_X)$) is injective; these homomorphisms therefore define a \emph{canonical injective homomorphism}
\[
\label{IV.20.1.4.1}
  i:\sh{O}_X\to\sh{M}_X
  \tag{20.1.4.1}
\]
which allows us to identify $\sh{O}_X$ with a sub-sheaf of $\sh{M}_X$.
Given a meromorphic function $\vphi\in M(X)$, we say that $\vphi$ is \emph{defined} in an open subset $U$ of $X$ if $\vphi|U$ is a \emph{section of $\sh{O}_U$} over $U$; the sheaf axioms show that, for a given section $\vphi$, there is a \emph{largest} open subset in which $\vphi$ is defined; it is called the \emph{domain of definition} of $\vphi$ and is denoted by $\operatorname{dom}(\vphi)$.
\end{env}

\begin{env}[20.1.5]\label{IV.20.1.5}
\EGAReaderTarget{ega4.c04.statement.20-1-5}{IV.20.1.5}
For every $\sh{O}_X$-module $\sh{F}$, we deduce from \sref{IV.20.1.4.1} a di-homomorphism consisting of $i$ and the homomorphism of sheaves of additive groups
\[
\label{IV.20.1.5.1}
  1_{\sh{F}}\otimes i:\sh{F}\to\sh{M}_X(\sh{F})=\sh{F}\otimes_{\sh{O}_X}\sh{M}_X.
  \tag{20.1.5.1}
\]
Note that the latter is no longer injective in general; when it is injective, we say that $\sh{F}$ is \emph{strictly torsion-free}: this means that for every open subset $U$ of $X$ and every section $s\in\Gamma(U,\sh{O}_X)$ that is a regular element of this ring, the homothety $z\mapsto sz$ of $\Gamma(U,\sh{F})$ is injective; this condition is evidently satisfied if $\sh{F}$ is \emph{locally free}.
\end{env}

\begin{proposition}[20.1.6]
\label{IV.20.1.6}
Let $X$ be a locally Noetherian prescheme and $\sh{F}$ a quasi-coherent $\sh{O}_X$-module.
For $\sh{F}$ to be strictly torsion-free, it is necessary and sufficient that $\operatorname{Ass}(\sh{F})\subset\operatorname{Ass}(\sh{O}_X)$.
\end{proposition}

\begin{proof}
Indeed, we are immediately reduced to the case where $X=\Spec(A)$ is affine and $\sh{F}=\widetilde M$, and we know that the elements $s$ of $A$ belonging to an ideal of $\operatorname{Ass}(M)$ are exactly those for which the homothety $z\mapsto sz$ is not injective (Bourbaki, \emph{Alg. comm.}, chap.~IV, \textsection1, n\textsuperscript{o}1, cor.~2 of prop.~2).
\end{proof}

\begin{env}[20.1.7]\label{IV.20.1.7}
\EGAReaderTarget{ega4.c04.statement.20-1-7}{IV.20.1.7}
If $u$ is a section of $\sh{M}_X(\sh{F})$ over $X$, we say that $u$ is \emph{defined} at a point $x\in X$ if there exists an open neighbourhood $V$ of $x$ in $X$ such that $u|V$ is the image of a section of $\sh{F}$ over $V$ under the di-homomorphism \sref{IV.20.1.5.1}.
We say that $u$ is \emph{defined} in an open subset $U$ of $X$ if it is defined at every point of $U$; there is again a largest open subset in which $u$ is defined, called the \emph{domain of definition} of $u$ and denoted by $\operatorname{dom}(u)$.
When $\sh{F}$ is strictly torsion-free, so that $\sh{F}$ is identified by \sref{IV.20.1.5.1} with a sub-sheaf of $\sh{M}_X(\sh{F})$, to say that $u$ is defined in $U$ means that $u|U$ is a \emph{section of $\sh{F}$ over $U$}.
\end{env}

\begin{env}[20.1.8]\label{IV.20.1.8}
\EGAReaderTarget{ega4.c04.statement.20-1-8}{IV.20.1.8}
In accordance with the general notation \sref[0\textsubscript{I}]{0.5.4.7}, we denote by $\sh{M}_X^*$ the sheaf of multiplicative groups such that $\Gamma(U,\sh{M}_X^*)$ is (for every open subset $U$ of $X$) the group of \emph{invertible elements} of $\Gamma(U,\sh{M}_X)$.
This sheaf is nothing other than the sheaf $\sh{S}(\sh{M}_X)$ defined in \sref{IV.20.1.3}: indeed, if $s\in\Gamma(U,\sh{S}(\sh{M}_X))$, then for every $x\in U$ there exists an open neighbourhood $V\subset U$ of $x$ such that $s|V$ is a regular element in the \emph{total ring of fractions} of $\Gamma(V,\sh{O}_X)$, and we know that such an element is necessarily invertible in this ring of fractions.
We will say that the sections of $\sh{M}_X^*$ over $X$ are the \emph{regular meromorphic functions} (note that we depart here from the terminology followed by certain authors, who call ``regular'' meromorphic functions those that are \emph{sections of $\sh{O}_X$}, identified with a sub-sheaf of $\sh{M}_X$).

Let $\sh{L}$ be an invertible $\sh{O}_X$-module \sref[0\textsubscript{I}]{0.5.4.1}; then it is clear that $\sh{M}_X(\sh{L})=\sh{L}\otimes_{\sh{O}_X}\sh{M}_X$
\oldpage[IV-4]{229}
is an invertible $\sh{M}_X$-module.
Let $U$ be an open subset such that $\sh{L}|U$ is isomorphic to $\sh{O}_U$; since every automorphism of $\sh{M}_U$ is multiplication by an invertible element of $\Gamma(U,\sh{M}_X)$ \sref[0\textsubscript{I}]{0.5.4.7}, it amounts to the same to say that a section $s\in\Gamma(U,\sh{M}_X(\sh{L}))$ has an invertible image in $\Gamma(U,\sh{M}_X)$ under \emph{one} isomorphism or under \emph{every} isomorphism onto $\Gamma(U,\sh{M}_X)$; in this case we will say that $s$ is a \emph{regular meromorphic section of $\sh{L}$ over $U$}; a section $s$ of $\sh{L}$ over $X$ will be called a \emph{regular meromorphic section of $\sh{L}$} if, for every open subset $U$ such that $\sh{L}|U$ is isomorphic to $\sh{O}_U$, $s|U$ is a regular meromorphic section of $\sh{L}$ over $U$.
We will denote by $(\sh{M}_X(\sh{L}))^*$ the sub-sheaf of $\sh{M}_X(\sh{L})$ such that, for every open subset $U$, $\Gamma(U,(\sh{M}_X(\sh{L}))^*)$ is the set of regular meromorphic sections of $\sh{L}$ over $U$.
Let $s$ be a meromorphic section of $\sh{L}$ over $X$ (i.e. a section of $\sh{M}_X(\sh{L})$); it defines a homomorphism $h_s:\sh{M}_X\to\sh{M}_X(\sh{L})$ which assigns $(s|U)t$ to every section $t$ of $\sh{M}_X$ over an open subset $U$.
It follows immediately from the above that for $s$ to be \emph{regular}, it is necessary and sufficient that $h_s$ be \emph{injective}, and in fact $h_s$ is then a \emph{bijective} homomorphism of $\sh{M}_X$ onto $\sh{M}_X(\sh{L})$, and its restriction to $\sh{M}_X^*$ is a bijection onto $(\sh{M}_X(\sh{L}))^*$.
We conclude that the homothety $t\mapsto ts$ is an isomorphism of $M(X)$ onto $M(X,\sh{L})$.
\end{env}

\begin{env}[20.1.9]
\label{IV.20.1.9}
Let $s$ be a regular meromorphic section of the invertible $\sh{O}_X$-module $\sh{L}$ over $X$; then, for every $\sh{O}_X$-module $\sh{F}$, $s$ likewise defines a homomorphism $h_s\otimes1_{\sh{F}}:\sh{M}_X(\sh{F})\to\sh{M}_X(\sh{F}\otimes_{\sh{O}_X}\sh{L})$, which is again \emph{bijective}.
\end{env}

\begin{env}[20.1.10]
\label{IV.20.1.10}
Let $s$ be a meromorphic section of an invertible $\sh{O}_X$-module $\sh{L}$ over $X$; for $s$ to be regular, it is necessary and sufficient that there exist a meromorphic section $s'$ of $\sh{L}^{-1}$ over $X$ such that the canonical image of $s\otimes s'$ in $\sh{M}_X$ \sref[0\textsubscript{I}]{0.5.4.3} is the unit section, and this section $s'$ is then unique: indeed, the necessity of the local existence of such a section is evident, and its local uniqueness implies its global existence (and uniqueness); moreover, the existence of $s'$ is trivially sufficient for $s$ to be regular.
We will set $s'=s^{-1}$.

Finally, if $\sh{L}'$ is a second invertible $\sh{O}_X$-module and $s$ (resp. $s'$) a regular meromorphic section of $\sh{L}$ (resp. $\sh{L}'$) over $X$, then $s\otimes s'$ is evidently a regular meromorphic section of $\sh{L}\otimes\sh{L}'$ over $X$.
\end{env}

\begin{env}[20.1.11]\label{IV.20.1.11}
\EGAReaderTarget{ega4.c04.statement.20-1-11}{IV.20.1.11}
If $f:X'\to X$ is a morphism of ringed spaces, there is in general no natural map associating to a meromorphic function on $X$ a meromorphic function on $X'$.
For example, if $X$ is the spectrum of an integral local ring $A$ and $X'$ that of its residue field $k$, there is no natural homomorphism from the field of fractions $K$ of $A$ to $k$, and one can associate an element of $k$ to an element of $K$ only if the latter is already in $A$.

In general, if $f=(\psi,\theta)$, let us denote, for every open subset $U$ of $X$, by $\sh{S}_f(U)$ the set of \emph{regular} sections $s\in\Gamma(U,\sh{O}_X)$ such that the image of $s$ under
\[
  \Gamma(\theta^\#):\Gamma(U,\sh{O}_X)\to\Gamma(f^{-1}(U),\sh{O}_{X'})
\]
is a \emph{regular} section.
It is immediate that $U\mapsto\sh{S}_f(U)$ is a \emph{sub-sheaf} of the sheaf of sets $\sh{S}(\sh{O}_X)$, which we denote by $\sh{S}_f$.
We set $\sh{M}_f=\sh{O}_X[\sh{S}_f^{-1}]$; it is a sub-sheaf
\oldpage[IV-4]{230}
of rings of $\sh{M}_X$, and from $\theta^\#:\psi^*(\sh{O}_X)\to\sh{O}_{X'}$ we canonically deduce a homomorphism of sheaves of rings $\theta^{\prime\#}:\psi^*(\sh{M}_f)\to\sh{M}_{X'}$ extending $\theta^\#$ (Bourbaki, \emph{Alg. comm.}, chap.~II, \textsection2, n\textsuperscript{o}1, prop.~2); whence, recalling that $f^*(\sh{M}_f)=\psi^*(\sh{M}_f)\otimes_{\psi^*(\sh{O}_X)}\sh{O}_{X'}$, a canonical homomorphism of $\sh{O}_{X'}$-algebras
\[
\label{IV.20.1.11.1}
  f^*(\sh{M}_f)\to\sh{M}_{X'}.
  \tag{20.1.11.1}
\]
For every meromorphic function $\vphi$ on $X$ that \emph{is a section of $\sh{M}_f$}, $\Gamma(\theta^{\prime\#})(\vphi)$ is a meromorphic function on $X'$, called the \emph{inverse image of $\vphi$ under $f$}, and denoted by $\vphi\circ f$ if this causes no confusion.

Likewise, if $\sh{F}$ is an $\sh{O}_X$-module, we set $\sh{M}_f(\sh{F})=\sh{F}\otimes_{\sh{O}_X}\sh{M}_f$, and from $\theta^{\prime\#}$ we immediately deduce a canonical homomorphism (which we also write $u\mapsto u\circ f$)
\[
  \Gamma(X,\sh{M}_f(\sh{F}))\to\Gamma(X',\sh{M}_{X'}(f^*(\sh{F}))).
\]
Moreover, if $u\in\Gamma(X,\sh{M}_f(\sh{F}))$ is defined \sref{IV.20.1.7} at a point $x$, then $u$ coincides, in a neighbourhood $U$ of $x$, with a section of the form $\sum_i h_i\otimes(t_i/s_i)$, where the $h_i$ belong to $\Gamma(U,\sh{F})$, the $t_i$ to $\Gamma(U,\sh{O}_X)$ and the $s_i$ to $\Gamma(U,\sh{S}_f)$.
Since by hypothesis the images of the $s_i$ in $\Gamma(f^{-1}(U),\sh{O}_{X'})$ are regular, we see that $u\circ f$ is defined at every point of $f^{-1}(U)$; in other words, we have
\[
\label{IV.20.1.11.2}
  f^{-1}(\operatorname{dom}(u))\subset\operatorname{dom}(u\circ f).
  \tag{20.1.11.2}
\]
We will see later \sref{IV.20.6.5}[(i)] examples (with $\sh{F}=\sh{O}_X$) where the two sides of \sref{IV.20.1.11.2} can be distinct.

Consider in particular the case where $\sh{M}_f=\sh{M}_X$; then, if $\sh{L}$ is an invertible $\sh{O}_X$-module, the image in $\sh{M}_{X'}(f^*(\sh{L}))$, under $\Gamma(\theta^{\prime\#})$, of a \emph{regular} meromorphic section of $\sh{L}$ over $X$ \sref{IV.20.1.8} is a \emph{regular} meromorphic section of $f^*(\sh{L})$ over $X'$, as follows immediately from the definition of these sections and from the fact that a ring homomorphism transforms an invertible element into an invertible element.

Let $f':X''\to X'$ be a second morphism of ringed spaces, and suppose that $\sh{M}_f=\sh{M}_X$ and $\sh{M}_{f'}=\sh{M}_{X'}$; then, if we set $f''=f\circ f'$, we also have $\sh{M}_{f''}=\sh{M}_X$, and we see immediately that for every meromorphic section $u$ of $\sh{F}$ over $X$, we have $u\circ f''=(u\circ f)\circ f'$.
\end{env}

\begin{proposition}[20.1.12]
\label{IV.20.1.12}
If the morphism $f:X'\to X$ is flat \sref[0\textsubscript{I}]{0.6.7.1}, then $\sh{M}_f=\sh{M}_X$, and the homomorphism $\vphi\mapsto\vphi\circ f$ is defined on the whole of $M(X)$.
Moreover, if $f$ is a (flat) morphism of ringed spaces in local rings, we have $\operatorname{dom}(\vphi\circ f)=f^{-1}(\operatorname{dom}(\vphi))$; if in addition $f$ is surjective (hence faithfully flat), the homomorphism $\vphi\mapsto\vphi\circ f$ is injective.
\end{proposition}

\begin{proof}
The first assertion follows from the fact that, if $B$ is an $A$-algebra that is a flat $A$-module, every element of $A$ that is not a divisor of $0$ in $A$ is not a divisor of $0$ in $B$ \sref[0\textsubscript{I}]{0.6.3.4}.
To prove the other assertions, note that, for every $x'\in X'$, if $x=f(x')$, then $\sh{O}_{X',x'}$ is a flat $\sh{O}_{X,x}$-module, and since the homomorphism $\sh{O}_{X,x}\to\sh{O}_{X',x'}$ is \emph{local} by hypothesis, it is injective (\sref[0\textsubscript{I}]{0.6.5.1} and \sref[0\textsubscript{I}]{0.6.6.2}); if we set $A=\sh{O}_{X,x}$, $B=\sh{O}_{X',x'}$, so that $A$ is identified with a subring of $B$, then $(f^*(\sh{M}_X))_{x'}$ is equal to $S^{-1}A\otimes_A B=S^{-1}B$, where $S$ is the set of regular elements of $A$, and $(\sh{M}_{X'})_{x'}$ is equal to $T^{-1}B$, where $T$ is the set
\oldpage[IV-4]{231}
of regular elements of $B$; and since we have seen that $S\subset T$, the homomorphism $S^{-1}B\to T^{-1}B$ is injective; in other words, this proves that the homomorphism \sref{IV.20.1.11.1} $f^*(\sh{M}_X)\to\sh{M}_{X'}$ is \emph{injective} (whence the last assertion of the statement).
The quotient $f^*(\sh{M}_X)/\sh{O}_{X'}$ is identified with a sub-$\sh{O}_{X'}$-module of $\sh{M}_{X'}/\sh{O}_{X'}$, and $(f^*(\sh{M}_X)/\sh{O}_{X'})_{x'}$ is identified with $(\sh{M}_X/\sh{O}_X)_x\otimes_{\sh{O}_{X,x}}\sh{O}_{X',x'}$.
Suppose then that $x\notin\operatorname{dom}(\vphi)$; the image of $\vphi_x$ in $(\sh{M}_X/\sh{O}_X)_x$ is then $\neq0$; by faithful flatness, we deduce that the same holds for the image of $(\vphi\circ f)_{x'}$ in $(\sh{M}_{X'}/\sh{O}_{X'})_{x'}$, hence $x'\notin\operatorname{dom}(\vphi\circ f)$, which completes the proof.
\end{proof}

\begin{remark}[20.1.13]
\label{IV.20.1.13}
Let $X$ be a \emph{reduced} complex analytic space; then the notion of meromorphic function on $X$ defined above coincides with the usual notion.
Consider, on the other hand, a prescheme $Y$ locally of finite type over the field $\bb{C}$; we know that one can then associate to $Y$ an analytic space $Y^\mathrm{an}$ having the same underlying topological space, and that the canonical morphism $f:Y^\mathrm{an}\to Y$ is \emph{flat}~\cite{IV-37}; by virtue of \sref{IV.20.1.12}, the canonical homomorphism $u\mapsto u\circ f$ from $M(Y)$ to $M(Y^\mathrm{an})$ is therefore everywhere defined and is injective; but it is not \emph{surjective} in general.
For example, when $Y=\bb{V}_{\bb{C}}^r$ (\textbf{Err}\textsubscript{\textbf{III}},~14) is the affine space of dimension $r$ over $\bb{C}$, $M(Y)$ is canonically identified with the field $R(Y)$ of rational functions on $Y$ \sref{IV.20.2.13}[(i)], whereas $M(Y^\mathrm{an})$ is the field of usual meromorphic functions on $\bb{C}^r$.
Because of this fact, it is often preferable, in Algebraic Geometry, to refrain from using the terminology introduced in this section, and to use the equivalent terminology of ``pseudo-function'', which will be defined below.
\end{remark}


\subsection{Pseudo-morphisms and pseudo-functions}
\EGAReaderTarget{ega4.c04.subsec.20-2}{IV.20.2}

The only ringed spaces considered in this section are \emph{preschemes}.

\begin{env}[20.2.1]\label{IV.20.2.1}
\EGAReaderTarget{ega4.c04.statement.20-2-1}{IV.20.2.1}
Recall \sref{IV.11.10.2} that, in a prescheme $X$, we say that an open subset $U$ is \emph{schematically dense} if, for every open subset $V$ of $X$, the canonical homomorphism $\Gamma(V,\sh{O}_X)\to\Gamma(V\cap U,\sh{O}_X)$ is \emph{injective}.

Consider two preschemes $X$, $Y$, and two \emph{schematically dense} open subsets $U$, $U'$ of $X$; we say that two morphisms $u:U\to Y$, $u':U'\to Y$ are \emph{equivalent} if there exists an open subset $U''\subset U\cap U'$, \emph{schematically dense} in $X$, such that $u|U''=u'|U''$.
Since it follows immediately from the definition of schematically dense open subsets that the intersection of two such open subsets is again one, it is immediate that the preceding relation is indeed an equivalence relation.
An equivalence class for this relation is called a \emph{pseudo-morphism} from $X$ to $Y$, or a \emph{strict rational map} from $X$ to $Y$.

If $S$ is a prescheme and $X$, $Y$ are two $S$-preschemes, we call \emph{pseudo-$S$-morphism} the equivalence class (for the preceding relation) of an $S$-morphism from a schematically dense open subset of $X$ to $Y$.
A pseudo-morphism is thus nothing other than a pseudo-$S$-morphism for $S=\Spec(\bb{Z})$.

We denote by $\operatorname{Ps.hom}(X,Y)$ (resp. $\operatorname{Ps.hom}_S(X,Y)$) the set of pseudo-morphisms (resp. of pseudo-$S$-morphisms) from $X$ to $Y$.
\end{env}

\begin{env}[20.2.2]\label{IV.20.2.2}
\EGAReaderTarget{ega4.c04.statement.20-2-2}{IV.20.2.2}
\oldpage[IV-4]{232}
It follows from the definition recalled above that if $U$ is an open subset schematically dense in $X$, then, for every open subset $V$ of $X$, $U\cap V$ is schematically dense in $V$.
If two morphisms $u:U\to Y$, $u':U'\to Y$ from schematically dense open subsets of $X$ to $Y$ are equivalent, it therefore follows that their restrictions $u|(V\cap U)$ and $u'|(V\cap U')$ are also equivalent morphisms (for the prescheme induced on $V$); their class is called the \emph{restriction} to $V$ of the pseudo-morphism $\omega$, the class of $u$, and this pseudo-morphism is denoted by $\omega|V$.
If $W\subset V$ is an open subset of $X$, it is clear that $(\omega|V)|W=\omega|W$.
This shows that the restriction maps define a \emph{presheaf} of sets $U\mapsto\operatorname{Ps.hom}(U,Y)$; note that this presheaf \emph{is not} a sheaf in general (cf. \sref{IV.20.2.16}); the associated sheaf is denoted by $\mathscr{P}\!s.\!\operatorname{hom}(X,Y)$.
For pseudo-$S$-morphisms, we see in the same way that $U\mapsto\operatorname{Ps.hom}_S(U,Y)$ is a presheaf of sets, whose associated sheaf we denote by $\mathscr{P}\!s.\!\operatorname{hom}_S(X,Y)$.

If $V$ is an open subset \emph{schematically dense} in $X$, then, for every open subset $U$ schematically dense in $X$, $U\cap V$ is also schematically dense in $X$; hence the map $\omega\mapsto\omega|V$ is a \emph{bijection} of the set $\operatorname{Ps.hom}(X,Y)$ (resp. $\operatorname{Ps.hom}_S(X,Y)$) onto $\operatorname{Ps.hom}(V,Y)$ (resp. $\operatorname{Ps.hom}_S(V,Y)$).
\end{env}

\begin{env}[20.2.3]\label{IV.20.2.3}
\EGAReaderTarget{ega4.c04.statement.20-2-3}{IV.20.2.3}
Given a pseudo-$S$-morphism $\omega$ from $X$ to $Y$, we say that it is \emph{defined} at a point $x\in X$ if there exist an open neighbourhood $V$ of $x$ in $X$, an open subset $U\subset V$ \emph{containing $x$ and schematically dense in $V$}, and finally an $S$-morphism $u:U\to Y$ whose class in $\operatorname{Ps.hom}_S(V,Y)$ is equal to $\omega|V$ \sref{IV.20.2.2}; we call \emph{domain of definition} of $\omega$, and denote by $\operatorname{dom}_S(\omega)$ (or simply $\operatorname{dom}(\omega)$ if $S=\Spec(\bb{Z})$), the set of $x\in X$ at which $\omega$ is defined; it is evidently an \emph{open subset} of $X$.
Moreover, for every open subset $W$ of $X$, we have
\[
\label{IV.20.2.3.1}
  \operatorname{dom}_S(\omega|W)=(\operatorname{dom}_S(\omega))\cap W
  \tag{20.2.3.1}
\]
by virtue of the property of schematically dense open subsets recalled in \sref{IV.20.2.2}.
\end{env}

\begin{proposition}[20.2.4]
\label{IV.20.2.4}
Suppose that $X$, $Y$ are $S$-preschemes such that the structural morphism $Y\to S$ is separated; then, if $\omega$ is a pseudo-$S$-morphism from $X$ to $Y$, $\operatorname{dom}_S(\omega)$ is the largest of the schematically dense open subsets $U$ of $X$ such that there is an $S$-morphism $u:U\to Y$ belonging to the class $\omega$.
\end{proposition}

\begin{proof}
It evidently suffices to prove that if $U$, $U'$ are two open subsets schematically dense in $X$ such that two $S$-morphisms $u:U\to Y$ and $u':U'\to Y$ are equivalent, then the restrictions of $u$ and $u'$ to $U\cap U'$ are equal.
Now, by hypothesis, there exists an open subset $U''\subset U\cap U'$, schematically dense in $X$, in which $u$ and $u'$ coincide; since $U''$ is also schematically dense in $U\cap U'$, our assertion follows from \sref{IV.11.10.1}[(d)].
\end{proof}

\begin{corollary}[20.2.5]
\label{IV.20.2.5}
Let $S$ be an $S_0$-scheme, and $X$, $Y$ two $S$-preschemes such that the composite morphism $Y\to S\to S_0$ is separated (which implies \sref[I]{I.5.5.1} that $Y\to S$ is separated too).
For every pseudo-$S$-morphism $\omega$ from $X$ to $Y$, we then have $\operatorname{dom}_S(\omega)=\operatorname{dom}_{S_0}(\omega)$.
In particular, if $Y$ is a scheme, we have $\operatorname{dom}_S(\omega)=\operatorname{dom}(\omega)$.
\end{corollary}

\begin{proof}
Indeed, it is clear that $\operatorname{dom}_S(\omega)\subset\operatorname{dom}_{S_0}(\omega)$ without any separatedness hypothesis; by virtue of \sref{IV.20.2.4}, it suffices to prove that if an $S_0$-morphism $u_0:U_0\to Y$ from a schematically dense open subset
\oldpage[IV-4]{233}
$U_0$ of $X$ to $Y$ is such that the composite $v_0:U_0\xrightarrow{u_0}Y\to S$ coincides with the restriction of the structural morphism $w_0:U_0\to S$ on an open subset $U\subset U_0$ schematically dense in $X$, then we have $v_0=w_0$.
But by virtue of the hypothesis that the morphism $S\to S_0$ is separated, this again follows from \sref{IV.11.10.1}[(d)].
\end{proof}

\begin{corollary}[20.2.6]
\label{IV.20.2.6}
Under the hypotheses of \sref{IV.20.2.4}, the presheaf $U\mapsto\operatorname{Ps.hom}_S(U,Y)$ is a sheaf.
\end{corollary}

\begin{proof}
Indeed, let $U$ be an open subset of $X$ and $(U_\alpha)$ a covering of $U$ by open subsets of $U$; suppose given, for each $\alpha$, a pseudo-$S$-morphism $\omega_\alpha$ from $U_\alpha$ to $Y$, in such a way that, for every pair of indices $\alpha$, $\beta$, the restrictions \sref{IV.20.2.2} of the pseudo-$S$-morphisms $\omega_\alpha$ and $\omega_\beta$ to $U_\alpha\cap U_\beta$ are equal; by virtue of \sref{IV.20.2.3.1}, this implies that $\operatorname{dom}_S(\omega_\alpha)\cap U_\beta=\operatorname{dom}_S(\omega_\beta)\cap U_\alpha$.
Then let $W$ be the union of the open subsets $\operatorname{dom}_S(\omega_\alpha)$, and, for every $\alpha$, let $u_\alpha$ be the unique $S$-morphism $\operatorname{dom}_S(\omega_\alpha)\to Y$ belonging to the class $\omega_\alpha$ \sref{IV.20.2.4}; because of the hypothesis and of \sref{IV.20.2.4}, the restrictions of $u_\alpha$ and $u_\beta$ to $\operatorname{dom}_S(\omega_\alpha)\cap\operatorname{dom}_S(\omega_\beta)$ are equal, so there exists a morphism $u:W\to Y$ whose restriction to each open subset $\operatorname{dom}_S(\omega_\alpha)$ is equal to $u_\alpha$; it is clear that $W$ is schematically dense in $U$, hence this morphism defines a pseudo-$S$-morphism $\omega$ from $U$ to $Y$ whose restrictions to the $U_\alpha$ are the $\omega_\alpha$.
\end{proof}

\begin{remark}[20.2.7]
\label{IV.20.2.7}
We know \sref{IV.11.10.4} that when $X$ is \emph{reduced}, it amounts to the same to say that an open subset of $X$ is dense in $X$ or that it is schematically dense in $X$; the notion of pseudo-morphism (resp. pseudo-$S$-morphism) from $X$ to $Y$ then coincides with that of \emph{rational map} (resp. \emph{rational $S$-map}) from $X$ to $Y$ \sref[I]{I.7.1.2}.
In the general case, the notion of pseudo-morphism seems more useful than that of rational map.
\end{remark}

\begin{env}[20.2.8]\label{IV.20.2.8}
\EGAReaderTarget{ega4.c04.statement.20-2-8}{IV.20.2.8}
We call \emph{pseudo-function} on $X$ a pseudo-morphism from $X$ to $\Spec(\bb{Z}[T])$ ($T$ an indeterminate), or, what amounts to the same, an $X$-pseudo-morphism from $X$ to $X\otimes_{\bb{Z}}\bb{Z}[T]$; it also amounts to the same \sref[I]{I.3.3.15} to say that a pseudo-function on $X$ is an \emph{equivalence class of sections of $\sh{O}_X$ over schematically dense open subsets of $X$}, two sections $g\in\Gamma(U,\sh{O}_X)$, $g'\in\Gamma(U',\sh{O}_X)$ over such open subsets being equivalent if there exists an open subset $U''\subset U\cap U'$, schematically dense in $X$, in which $g$ and $g'$ coincide.
We can apply \sref{IV.20.2.4} here with $S=X$ and $Y=X\otimes_{\bb{Z}}\bb{Z}[T]$; hence, for every pseudo-function $\vphi$ on $X$, there exists a \emph{largest} schematically dense open subset $\operatorname{dom}(\vphi)$ of $X$ such that there is a section of $\sh{O}_X$ over $\operatorname{dom}(\vphi)$ belonging to the class $\vphi$.
It is clear that the sheaf $\mathscr{P}\!s.\!\operatorname{hom}_X(X,X\otimes_{\bb{Z}}\bb{Z}[T])$ is here a \emph{sheaf of rings}, and even an $\sh{O}_X$-\emph{algebra}, which we will denote by $\sh{M}'_X$; it follows from \sref{IV.20.2.6} that, for every open subset $U$ of $X$, $\Gamma(U,\sh{M}'_X)$ is equal to the set of pseudo-morphisms from $U$ to $\Spec(\bb{Z}[T])$; we set $M'(X)=\Gamma(X,\sh{M}'_X)$.
To say that a section $\vphi$ of $\sh{M}'_X$ over $U$ is \emph{invertible} in the ring $\Gamma(U,\sh{M}'_X)$ means that there exists an open subset $U'$ schematically dense in $\operatorname{dom}(\vphi)$, hence in $U$, such that, if $g$ is the unique section of $\sh{O}_X$ over $\operatorname{dom}(\vphi)$ belonging to $\vphi$, then $g|U'$ is \emph{invertible} in $\Gamma(U',\sh{O}_X)$.
It follows from \sref[I]{I.3.3.15} that, in the canonical correspondence between $\Gamma(V,\sh{O}_X)$ and $\Hom(V,\bb{Z}[T])$ ($V$ an open subset of $X$), the invertible elements of $\Gamma(V,\sh{O}_X)$ correspond to the morphisms that
\oldpage[IV-4]{234}
\emph{factor} as $V\to\Spec(\bb{Z}[T,T^{-1}])\to\Spec(\bb{Z}[T])$.
We conclude that the sheaf ${\sh{M}'}^*_X$ of germs of invertible sections of $\sh{M}'_X$ is canonically identified with the sheaf $\mathscr{P}\!s.\!\operatorname{hom}_X(X,X\otimes_{\bb{Z}}\bb{Z}[T,T^{-1}])$.
\end{env}

\begin{lemma}[20.2.9]
\label{IV.20.2.9}
Let $U$ be an open subset of $X$ and $s$ a regular element of the ring $\Gamma(U,\sh{O}_X)$; then the open set $U_s$ of points $x\in U$ such that $s(x)\neq0$ is schematically dense in $U$.
\end{lemma}

\begin{proof}
Indeed, let $V$ be an open subset of $U$ and $t$ a section of $\sh{O}_X$ over $V$ such that $t|(V\cap U_s)=0$.
For every affine open subset $W\subset V$, there then exists an integer $n$ such that $s^nt|W=0$ \sref[I]{I.1.4.1}; but since $s$ is a \emph{regular} element of $\Gamma(U,\sh{O}_X)$, this implies $t|W=0$, whence $t=0$.
\end{proof}

\begin{env}[20.2.10]
\label{IV.20.2.10}
Consider a meromorphic function $f\in M(X)$ \sref{IV.20.1.4}; then $\operatorname{dom}(f)$ is \emph{schematically dense} in $X$: indeed, every point of $X$ has an open neighbourhood $U$ such that there exists a section $s\in\Gamma(U,\sh{O}_X)$ that is a regular element of this ring and for which $s(f|U)\in\Gamma(U,\sh{O}_X)$; since $s|U_s$ is invertible, we conclude that $f|U_s\in\Gamma(U_s,\sh{O}_X)$, hence $U_s\subset\operatorname{dom}(f)$ by definition \sref{IV.20.1.4}, and our assertion therefore follows from Lemma \sref{IV.20.2.9}.
We can therefore associate to $f$ the pseudo-function equal to the class of the section of $\sh{O}_X$ over $\operatorname{dom}(f)$ that is the restriction of $f$; proceeding in the same way with $X$ replaced by an arbitrary open subset of $X$, we thus obtain a \emph{canonical homomorphism of $\sh{O}_X$-algebras}
\[
\label{IV.20.2.10.1}
  \sh{M}_X\to\sh{M}'_X
  \tag{20.2.10.1}
\]
which, by restriction, evidently gives a homomorphism of sheaves of abelian groups
\[
\label{IV.20.2.10.2}
  \sh{M}_X^*\to{\sh{M}'}^*_X
  \tag{20.2.10.2}
\]
for the sheaves of invertible germs of sections of $\sh{M}_X$ and $\sh{M}'_X$.
\end{env}

\begin{proposition}[20.2.11]\label{IV.20.2.11}
\EGAReaderTarget{ega4.c04.prop.20-2-11}{IV.20.2.11}
\medskip
\begin{enumerate}[label=\emph{(\roman*)}]
  \item The canonical homomorphism \emph{\sref{IV.20.2.10.1}} (and consequently also \emph{\sref{IV.20.2.10.2}}) is \emph{injective}.
  \item Suppose either that $X$ is locally Noetherian, or that $X$ is reduced and the set of its irreducible components is locally finite.
  Then the canonical homomorphism \emph{\sref{IV.20.2.10.1}} (and consequently also \emph{\sref{IV.20.2.10.2}}) is \emph{bijective}.
\end{enumerate}
\end{proposition}

\begin{proof}
The questions under consideration being local on $X$, we may restrict to considering the case where $X=\Spec(A)$ is affine, and to showing then that the canonical homomorphism $M(X)\to M'(X)$ is always injective, and that it is bijective in the cases considered in~(ii).
Taking into account the definition of the sheaf $\sh{M}_X$ \sref{IV.20.1.3}, we may moreover remark that \sref{IV.20.2.10.1} in fact comes from a homomorphism of \emph{presheaves}
\[
  \Gamma(U,\sh{O}_X)[\Gamma(U,\sh{S})^{-1}]\to\Gamma(U,\sh{M}'_X)
\]
and it suffices to show that, for $U$ affine, the latter is injective (resp. bijective under the hypotheses of~(ii)).
Denoting by $S$ the set of regular elements of $A$ (so that $S^{-1}A$ is \emph{the total ring of fractions} of $A$), we must therefore consider the canonical homomorphism
\[
\label{IV.20.2.11.1}
  S^{-1}A\to M'(X)
  \tag{20.2.11.1}
\]
\oldpage[IV-4]{235}
and show that it is injective (resp. bijective under the conditions of~(ii)).
Now, we can write $S^{-1}A=\varinjlim A_t$, where $t$ runs through the set of \emph{regular} elements of $A$, ordered by the relation ``$t$ is a divisor of $t'$'' \sref[0]{0.1.4.5}; we can also write $A_t=\Gamma(\mathrm{D}(t),\sh{O}_X)$.
On the other hand, by definition we have $M'(X)=\varinjlim\Gamma(U,\sh{O}_X)$, where $U$ runs through the set of schematically dense open subsets of $X$ (ordered by the relation $\supset$), and the map \sref{IV.20.2.11.1} is nothing other than the canonical map coming from the fact that the $\mathrm{D}(t)$ constitute a subset of the set of open subsets schematically dense in $X$ \sref{IV.20.2.9}.
Note that, by definition of a schematically dense open subset, the homomorphism $\Gamma(U,\sh{O}_X)\to\Gamma(U',\sh{O}_X)$ for two such open subsets $U'\subset U$ is always \emph{injective}, so the same is true of the homomorphisms $\Gamma(U,\sh{O}_X)\to M'(X)$, and we conclude that \sref{IV.20.2.11.1} is injective.
To prove that \sref{IV.20.2.11.1} is bijective, it suffices to show that the set of the $\mathrm{D}(t)$ is \emph{cofinal} in the set of schematically dense open subsets, in other words that for such an open subset $U$, there exists $t$ regular in $A$ such that $\mathrm{D}(t)\subset U$.
Now, when $X$ is reduced and the irreducible components of $X$ form a locally finite set, this set is finite since $X$ has been supposed affine; in other words, $A$ has only finitely many minimal prime ideals $\mathfrak{p}_i$; since $A$ is reduced, the intersection of the $\mathfrak{p}_i$ is reduced to $0$, and to say that $t$ is regular is therefore equivalent to saying that $t$ belongs to none of the $\mathfrak{p}_i$; we therefore conclude by the reasoning of \sref[I]{I.7.1.9.1}.
When $A$ is Noetherian, to say that $U=X-Y$ (where $Y=V(\mathfrak{J})$ is closed in $X$) is schematically dense means \sref{IV.5.10.2} that $Y$ \emph{does not meet} $\operatorname{Ass}(\sh{O}_X)$, and by virtue of Bourbaki, \emph{Alg.~comm.}, chap.~IV, \textsection1, n\textsuperscript{o}~4, prop.~8, this implies the existence of a $t\in\mathfrak{J}$ such that $t$ is $A$-regular, hence $U\supset\mathrm{D}(t)$.

In the course of this proof we have moreover proved the following
\end{proof}

\begin{corollary}[20.2.12]
\label{IV.20.2.12}
If $X=\Spec(A)$, where $A$ is Noetherian, or reduced and having only finitely many minimal prime ideals, the schematically dense open subsets of $X$ are those that contain an open subset of the form $\mathrm{D}(t)$, where $t$ is regular in $A$, and $M(X)=M'(X)$ is the total ring of fractions $S^{-1}A$, where $S$ is the set of regular elements of $X$.
\end{corollary}

\begin{remarks}[20.2.13]\label{IV.20.2.13}
\EGAReaderTarget{ega4.c04.statement.20-2-13}{IV.20.2.13}
\begin{enumerate}[label=(\roman*)]
  \item Let $\vphi$ be an element of $M(X)$ and $\vphi'$ its image in $M'(X)$; we evidently have, by definition (\sref{IV.20.1.4} and \sref{IV.20.2.8}), $\operatorname{dom}(\vphi)\subset\operatorname{dom}(\vphi')$; but in fact we have the equality $\operatorname{dom}(\vphi)=\operatorname{dom}(\vphi')$, for there exists a section of $\sh{O}_X$ over $\operatorname{dom}(\vphi')$ belonging to the class $\vphi'$, and the corresponding meromorphic function is equal to $\vphi$ \sref{IV.20.2.11}[(i)], hence $\operatorname{dom}(\vphi')\subset\operatorname{dom}(\vphi)$.

  \item We have already noted that when $X$ is \emph{reduced}, we have $\sh{M}'_X=\sh{R}(X)$ (sheaf of rational functions on $X$ \sref[I]{I.7.3.2}); if moreover the irreducible components of $X$ form a locally finite set, we have $\sh{M}_X=\sh{M}'_X=\sh{R}(X)$.
  In general, since every schematically dense open set is dense, we have a canonical homomorphism $\sh{M}'_X\to\sh{R}(X)$, but even when $X$ is locally Noetherian, this homomorphism is not necessarily injective.
  For example, if $X=\Spec(A)$, where $A$ is a Noetherian ring such that $\operatorname{Ass}(A)$ contains embedded prime ideals (which implies that $A$ is not reduced), we have seen that $M(X)$ and $M'(X)$ are identified with the total ring of fractions $S^{-1}A$,
  \oldpage[IV-4]{236}
  where $S$ is the set of regular elements of $A$, the complement of the union of the ideals $\mathfrak{p}\in\operatorname{Ass}(A)$; on the other hand, $R(X)$ is identified with $Q^{-1}A$, where $Q$ is the complement of the union of the \emph{minimal} prime ideals of $A$ \sref[I]{I.7.1.9}, and the canonical homomorphism $A\to Q^{-1}A$ (and \emph{a fortiori} $S^{-1}A\to Q^{-1}A$) is therefore not injective, since there exist in $A-Q$ elements $\neq0$ of $A$ annihilated by elements of $Q$ (Bourbaki, \emph{Alg.~comm.}, chap.~IV, \textsection1, n\textsuperscript{o}~1, cor.~2 of prop.~1).

  \item Note that even when $X$ is locally Noetherian, the $\sh{O}_X$-module $\sh{M}'_X$ \emph{is not necessarily quasi-coherent}.
  Consider for example a Noetherian local ring $A$ of dimension $\geq2$ whose maximal ideal $\mathfrak{m}$ is such that $\mathfrak{m}\in\operatorname{Ass}(A)$ and such that, if we set $X=\Spec(A)$, the scheme induced on the open subset $U$ of $X$ complementary to $\{\mathfrak{m}\}$ is \emph{integral}.
  The only regular elements of $A$ are then the invertible elements, so $\Gamma(X,\sh{M}'_X)=M'(X)=A$; if $\sh{M}'_X$ were quasi-coherent, it would therefore be equal to $\sh{O}_X$; but since $U$ is an integral scheme, $\sh{M}'_X|U=R(U)$ is a \emph{simple} sheaf \sref[I]{I.7.3.5}, whereas $\sh{O}_U$ is not a simple sheaf since $\dim(A)\geq2$.

  It remains to give an example of a ring $A$ having the preceding properties.
  It suffices to consider an integral local ring $B$ of dimension $\geq2$ with residue field $k$, and to take $A=B\oplus k$ with the multiplication law $(b,x)(b',x')=(bb',bx'+b'x)$.

  \item If $X$ is locally Noetherian, it follows from \sref{IV.5.10.2} that the schematically dense open subsets of $X$ are those that \emph{contain} the set $\operatorname{Ass}(\sh{O}_X)$.
\end{enumerate}
\end{remarks}

\begin{env}[20.2.14]\label{IV.20.2.14}
\EGAReaderTarget{ega4.c04.statement.20-2-14}{IV.20.2.14}
Let $X$ be a prescheme and $\sh{F}$ a quasi-coherent and \emph{strictly torsion-free} $\sh{O}_X$-module \sref{IV.20.1.5}, so that $\sh{F}$ is identified with a sub-$\sh{O}_X$-module of $\sh{M}_X(\sh{F})$.
For every meromorphic section $u$ of $\sh{F}$ over $X$, we call \emph{ideal of denominators of $u$} the annihilator $\sh{J}$ of the section $\overline{u}$, the image of $u$ in $\sh{M}_X(\sh{F})/\sh{F}$.
The ideal $\sh{J}$ is \emph{quasi-coherent}: indeed, the question being local on $X$, we may restrict to the case where $X$ is affine, and there exists a section $s\in\Gamma(X,\sh{S}(\sh{O}_X))$ such that $v=su\in\Gamma(X,\sh{F})$.
To say that, for an open subset $U\subset X$, a section $f\in\Gamma(U,\sh{O}_X)$ belongs to $\Gamma(U,\sh{J})$ means that $f(u|U)\in\Gamma(U,\sh{F})$, and since $s|U$ is a regular element of $\Gamma(U,\sh{O}_X)$ and $\sh{F}$ is strictly torsion-free, the preceding relation is also equivalent to $f((su)|U)\in\Gamma(U,s\sh{F})$; if $\overline{v}$ is the section of $\sh{F}/s\sh{F}$ that is the canonical image of $v$, we therefore see that $\sh{J}$ is the kernel of the homomorphism $\sh{O}_X\to\sh{F}/s\sh{F}$ obtained by multiplication by the section $\overline{v}$.
Since $\sh{F}/s\sh{F}$ is quasi-coherent, so is $\sh{J}$.

It follows immediately from the preceding definition that $\operatorname{dom}(u)$ is \emph{the open complement of the closed sub-prescheme of $X$ defined by the ideal of denominators of $u$}.
\end{env}

\begin{proposition}[20.2.15]
\label{IV.20.2.15}
Let $f:X'\to X$ be a morphism, $\sh{F}$ a quasi-coherent $\sh{O}_X$-module, and $\vphi$ a section of $\sh{M}_f(\sh{F})$ over $X$ \sref{IV.20.1.11}.
Then $f^{-1}(\operatorname{dom}(\vphi))$ is schematically dense in $X'$.
\end{proposition}

\begin{proof}
The question being evidently local on $X$ and $X'$, we may suppose $X=\Spec(A)$, $X'=\Spec(A')$ affine, $\sh{F}=\widetilde{M}$, and $\vphi=h\otimes(1/s)$, where $h\in M$ and $s$ is a regular element of $A$ whose image $s'$ in $A'$ is a regular element.
We then know \sref{IV.20.2.9} that $\mathrm{D}(s')$ is a schematically dense open subset of $X'$, and it is the inverse image under $f$ of $\mathrm{D}(s)$, which is contained in $\operatorname{dom}(\vphi)$.
\end{proof}

\begin{remark}[20.2.16]
\label{IV.20.2.16}
When $Y$ is not separated, the presheaf $U\mapsto\operatorname{Ps.hom}_S(U,Y)$ on $X$ is not necessarily a sheaf, even when $X$ is Noetherian, as the following example shows.
We will take for $X$ the spectrum of a Noetherian semi-local ring $A$ of dimension $\geq2$, having exactly two maximal
\oldpage[IV-4]{237}
ideals $\mathfrak{m}'$, $\mathfrak{m}''$ (so that $X$ has exactly two closed points $x'$, $x''$), in such a way that $\mathfrak{m}'$ and $\mathfrak{m}''$ belong to $\operatorname{Ass}(A)$ and the open subset $U=X-\{x',x''\}$ is \emph{integral}.
Let $U'=X-\{x'\}$, $U''=X-\{x''\}$, so that $U=U'\cap U''$.
Note that the only schematically dense open subsets of $X$ are those that contain $x'$ and $x''$ \sref{IV.20.2.13}[(iv)], so they are necessarily identical to $X$.
To obtain a counterexample, it will therefore suffice to define two $S$-morphisms $u':U'\to Y$, $u'':U''\to Y$ (with $S=X$) whose restrictions to $U$ belong to the same pseudo-morphism from $U$ to $Y$, but which are such that, for every neighbourhood $V'$ of $x''$ in $U'$ and every neighbourhood $V''$ of $x'$ in $U''$, the restrictions of $u'$ and $u''$ to $V'\cap V''$ are distinct.
For this, consider an irreducible closed subset $Z$ of $X$ containing $x'$ and $x''$, distinct from $X$; let $Y$ be the $X$-prescheme obtained by gluing two preschemes $Y'$, $Y''$ isomorphic to $X$ along the everywhere dense open subset $X-Z$ \sref[I]{I.2.3.1}.
We then take for $u'$ and $u''$ the respective restrictions to $U'$ and $U''$ of the isomorphisms inverse to the structural isomorphisms $Y'\to X$, $Y''\to X$.
Since $V'$ and $V''$ contain the generic point of $Z$, the restrictions of $u'$ and $u''$ to $V'\cap V''$ are distinct, but the restrictions of $u'$ and $u''$ to $U-(U\cap Z)$ are equal, and $U-(U\cap Z)$ is a dense open subset of $U$, hence schematically dense since $U$ is reduced.

It therefore only remains to define $A$ and $Z$ so as to have the preceding properties.
Let $X_0=\Spec(B)$ be an integral affine scheme (for example the affine plane over a field $k$), and $Y$ an irreducible closed subset of $X_0$ (for example an affine line) containing two distinct closed points $x'$ and $x''$, corresponding to the maximal ideals $\mathfrak{n}'$, $\mathfrak{n}''$ of $B$.
Consider the ring $C=B\oplus(B/\mathfrak{n}'\oplus B/\mathfrak{n}'')$ with the multiplication $(b,z)(b',z')=(bb',bz'+b'z)$.
If $X_1=\Spec(C)$, we have $X_0=(X_1)_\mathrm{red}$, and $X_1$ is reduced except at the points $x'$, $x''$.
It then suffices to take $A=R^{-1}C$, where $R$ is the complement of the union of the maximal ideals of $C$ at the points $x'$, $x''$, and for $Z$ the trace of $Y$ on $X=\Spec(A)$.
\end{remark}


\subsection{Composition of pseudo-morphisms}
\EGAReaderTarget{ega4.c04.subsec.20-3}{IV.20.3}

\begin{env}[20.3.1]\label{IV.20.3.1}
\EGAReaderTarget{ega4.c04.statement.20-3-1}{IV.20.3.1}
Let $X$, $Y$, $Z$ be three preschemes, $\omega$ a pseudo-morphism from $X$ to $Y$, $f:Y\to Z$ a morphism.
It is clear that if $U'$, $U''$ are two schematically dense open subsets of $X$ and $u':U'\to Y$, $u'':U''\to Y$ two morphisms of the class $\omega$, then the morphisms $f\circ u'$ and $f\circ u''$ are equivalent (for the relation defined in \sref{IV.20.2.1}); hence, for all the morphisms $u$ of the class $\omega$, the morphisms $f\circ u$ belong to one and the same equivalence class, which we denote by $f\circ\omega$ and call the \emph{composite of $f$ and $\omega$}; it is a pseudo-morphism from $X$ to $Z$.
We have $\operatorname{dom}(f\circ\omega)=\operatorname{dom}(\omega)$.
If $g:Z\to T$ is a morphism, it is clear that $g\circ(f\circ\omega)=(g\circ f)\circ\omega$.
\end{env}

\begin{env}[20.3.2]\label{IV.20.3.2}
\EGAReaderTarget{ega4.c04.statement.20-3-2}{IV.20.3.2}
Now suppose given a pseudo-$S$-morphism $\omega$ from $X$ to $Y$, where $Y$ is \emph{separated over $S$}, so that there exists an $S$-morphism $u:\operatorname{dom}_S(\omega)\to Y$ of the class $\omega$ \sref{IV.20.2.4}.
Let $f:X'\to X$ be an $S$-morphism such that the open subset $V'=f^{-1}(\operatorname{dom}_S(\omega))$ is schematically dense in $X'$; then we say that the class (for the equivalence relation of \sref{IV.20.2.1}) of the $S$-morphism $u\circ(f|U')$ (a class which is defined by virtue of the preceding hypothesis) is the \emph{composite of $\omega$ and $f$}, a pseudo-$S$-morphism which we denote by $\omega\circ f$; it is clear that we have
\[
\label{IV.20.3.2.1}
  f^{-1}(\operatorname{dom}_S(\omega))\subset\operatorname{dom}_S(\omega\circ f)
  \tag{20.3.2.1}
\]

For a given $f:X'\to X$, we denote by $\operatorname{Ps.hom}_S(X,Y)^f$ the set of pseudo-$S$-morphisms $\omega$ from $X$ to $Y$ that satisfy the preceding condition.
If $\omega$ is such a pseudo-$S$-morphism, it is clear that for every open subset $V$ of $X$,
\[
  f^{-1}(\operatorname{dom}_S(\omega|V))=f^{-1}(V\cap\operatorname{dom}_S(\omega))=f^{-1}(V)\cap f^{-1}(\operatorname{dom}_S(\omega))
\]
is schematically dense in $f^{-1}(V)$; hence, if $f^V:f^{-1}(V)\to V$ is the restriction of $f$, the composite $(\omega|V)\circ f^V$ is defined and equal to $(\omega\circ f)|f^{-1}(V)$.
We thus define a
\oldpage[IV-4]{238}
canonical restriction map $\operatorname{Ps.hom}_S(X,Y)^f\to\operatorname{Ps.hom}_S(V,Y)^{f^V}$, which allows us to define a \emph{presheaf} of sets $V\mapsto\operatorname{Ps.hom}_S(V,Y)^{f^V}$ on $X$, a subpresheaf of $V\mapsto\operatorname{Ps.hom}_S(V,Y)$, whence an associated sheaf of sets, which we denote by $\mathscr{P}\!s.\!\operatorname{hom}_S(X,Y)^f$.
Moreover, for every open subset $V$ of $X$, we have a map $\omega\mapsto\omega\circ f^V$ from $\operatorname{Ps.hom}_S(V,Y)^{f^V}$ to $\operatorname{Ps.hom}_S(f^{-1}(V),Y)$, which therefore defines an $f$-morphism of sheaves of sets
\[
  \mathscr{P}\!s.\!\operatorname{hom}_S(X,Y)^f\to\mathscr{P}\!s.\!\operatorname{hom}_S(X',Y).
\]
\end{env}

\begin{env}[20.3.3]
\label{IV.20.3.3}
Now let $f':X''\to X'$ be an $S$-morphism such that the open subset $f'^{-1}(f^{-1}(\operatorname{dom}_S(\omega)))$ is schematically dense in $X''$; then $\omega\circ(f\circ f')$ is defined and $u\circ f\circ f'$ belongs to this pseudo-$S$-morphism; moreover, by virtue of \sref{IV.20.3.2.1}, $f'^{-1}(\operatorname{dom}_S(\omega\circ f))$ is \emph{a fortiori} schematically dense in $X''$, so $(\omega\circ f)\circ f'$ is also defined and $u\circ f\circ f'$ belongs to this pseudo-$S$-morphism; hence we have $(\omega\circ f)\circ f'=\omega\circ(f\circ f')$.
On the other hand, for every $S$-morphism $g:Y\to Z$, we have $\operatorname{dom}_S(g\circ\omega)=\operatorname{dom}_S(\omega)$ \sref{IV.20.3.1}, so $(g\circ\omega)\circ f$ is defined, and $g\circ u\circ f$ belongs to this pseudo-$S$-morphism, which shows that $(g\circ\omega)\circ f=g\circ(\omega\circ f)$.
\end{env}

\begin{env}[20.3.4]
\label{IV.20.3.4}
The most important case in definition \sref{IV.20.3.2} is the one where $\mathscr{P}\!s.\!\operatorname{hom}_S(X,Y)^f=\mathscr{P}\!s.\!\operatorname{hom}_S(X,Y)$: for this it suffices that, for every open subset $U$ of $X$ and every open subset $V$ schematically dense in $U$, $f^{-1}(V)$ be schematically dense in $f^{-1}(U)$; when this is so, then, for \emph{every} open subset $U$ of $X$ and \emph{every} pseudo-$S$-morphism $\omega:U\to Y$, we can define the composite $\omega\circ f^U$, \emph{even if $Y$ is not separated over $S$}.
Indeed, if $u':U'\to Y$ and $u'':U''\to Y$ are two morphisms of the class $\omega$, they coincide on an open subset $U_0$ schematically dense in $U$, so the composite morphisms $f^{-1}(U')\to U'\xrightarrow{u'}Y$ and $f^{-1}(U'')\to U''\xrightarrow{u''}Y$ coincide on $f^{-1}(U_0)$, which by hypothesis is schematically dense in $f^{-1}(U)$; we can therefore define $\omega\circ f^U$ as the class of any one of the morphisms $f^{-1}(U')\to U'\xrightarrow{u'}Y$, where $u'$ belongs to $\omega$.
\end{env}

\begin{proposition}[20.3.5]\label{IV.20.3.5}
\EGAReaderTarget{ega4.c04.prop.20-3-5}{IV.20.3.5}
Let $X$, $X'$ be two preschemes, $f:X'\to X$ a morphism.
Then, in each of the following three cases, for every open subset $U$ of $X$ and every open subset $V$ schematically dense in $U$, $f^{-1}(V)$ is schematically dense in $f^{-1}(U)$:
\begin{enumerate}[label=\emph{(\roman*)}]
  \item $f$ is flat and locally of finite presentation.
  \item $f$ is flat and the underlying space of $X$ is locally Noetherian.
  \item $X'$ is reduced, the set of irreducible components of $X$ is locally finite, and every irreducible component of $X'$ dominates an irreducible component of $X$.
\end{enumerate}
\end{proposition}

\begin{proof}
Assertion~(i) follows from \sref{IV.11.10.5}[(ii), (b)]; assertion~(ii) follows from \sref{IV.11.10.5}[(ii), (a)], because then every open subset $V$ in $U$ is retrocompact, in other words the canonical injection $j:V\to U$ is a quasi-compact morphism.
Finally, to prove~(iii), note that since $X'$ is reduced, it suffices to show that for every open subset $U$ of $X$ and every open subset $V$ dense in $U$, $f^{-1}(V)$ is \emph{dense} in $f^{-1}(U)$.
Now, for $f^{-1}(V)$ to be dense in $f^{-1}(U)$, it suffices that $f^{-1}(V)$ contain all the maximal points of $X'$ belonging to $f^{-1}(U)$; the conclusion therefore follows from the hypothesis that the image under $f$ of every
\oldpage[IV-4]{239}
maximal point of $X'$ belonging to $f^{-1}(U)$ is a maximal point of $X$ belonging to $U$, hence to $V$, since the set of irreducible components of $X$ is locally finite.
\end{proof}

\begin{env}[20.3.6]\label{IV.20.3.6}
\EGAReaderTarget{ega4.c04.statement.20-3-6}{IV.20.3.6}
Let $X$, $Y$ be two $S$-preschemes; suppose that $X$ satisfies one of the following two hypotheses:
\begin{enumerate}[label=\emph{\alph*)}]
  \item $X$ is locally Noetherian;
  \item $X$ is reduced and the set of its irreducible components is locally finite.
\end{enumerate}
Then, for every $x\in X$, the canonical $S$-morphism $j:\Spec(\sh{O}_{X,x})\to X$ is flat and satisfies condition~(ii) of \sref{IV.20.3.5} in case~a), and condition~(iii) of \sref{IV.20.3.5} in case~b).
For every pseudo-$S$-morphism $\omega$ from $X$ to $Y$, the composite $\omega\circ j$ is therefore defined and is a pseudo-$S$-morphism from $\Spec(\sh{O}_{X,x})$ to $Y$, called the \emph{restriction} of $\omega$ to $\Spec(\sh{O}_{X,x})$.
Note now that if $X$ satisfies condition~a) (resp. b)) of \sref{IV.20.3.6}, then so does every prescheme induced on an open subset $U$ of $X$ containing $x$.
By passing to the inductive limit, we therefore deduce, from the canonical maps $\operatorname{Ps.hom}_S(U,Y)\to\operatorname{Ps.hom}_S(\Spec(\sh{O}_{X,x}),Y)$ thus obtained, a canonical map
\[
\label{IV.20.3.6.1}
  (\mathscr{P}\!s.\!\operatorname{hom}_S(X,Y))_x\to\operatorname{Ps.hom}_S(\Spec(\sh{O}_{X,x}),Y)
  \tag{20.3.6.1}
\]
where the left-hand side is the stalk at the point $x$ of the sheaf $\mathscr{P}\!s.\!\operatorname{hom}_S(X,Y)$, the set of germs at $x$ of pseudo-$S$-morphisms from open neighbourhoods of $x$ to $Y$.
\end{env}

\begin{proposition}[20.3.7]
\label{IV.20.3.7}
Under the hypotheses of \sref{IV.20.3.6}, the canonical map \sref{IV.20.3.6.1} is injective.
If moreover $Y$ is locally of finite presentation over $S$, the map \sref{IV.20.3.6.1} is bijective.
\end{proposition}

\begin{proof}
By applying the method of \sref{IV.8.1.2}[(a)], this proposition will be a particular case of the following one:
\end{proof}

\begin{proposition}[20.3.8]\label{IV.20.3.8}
\EGAReaderTarget{ega4.c04.prop.20-3-8}{IV.20.3.8}
With the notation of \sref{IV.8.8.1}, suppose $X_\alpha$ quasi-compact (resp. quasi-compact and quasi-separated) and $Y_\alpha$ locally of finite type (resp. locally of finite presentation) over $S_\alpha$.
Suppose moreover that one of the following conditions is satisfied:
\begin{enumerate}[label=\emph{(\roman*)}]
  \item The transition morphisms $S_\mu\to S_\lambda$ (for $\lambda\leq\mu$) are flat, and the $X_\lambda$ and $X$ are Noetherian.
  \item The $X_\lambda$ are reduced, the set of irreducible components of $X$ and of each of the $X_\lambda$ is finite, and, for $\lambda\leq\mu$, every irreducible component of $X_\mu$ dominates an irreducible component of $X_\lambda$.
\end{enumerate}
Then the canonical map
\[
\label{IV.20.3.8.1}
  \varinjlim\operatorname{Ps.hom}_{S_\lambda}(X_\lambda,Y_\lambda)\to\operatorname{Ps.hom}_S(X,Y)
  \tag{20.3.8.1}
\]
is injective (resp. bijective).
\end{proposition}

\begin{proof}
Note first that, in case~(i), the morphisms $X_\mu\to X_\lambda$ (for $\lambda\leq\mu$) and $X\to X_\lambda$ are flat, so it follows from \sref{IV.20.3.4} and \sref{IV.20.3.5} that the canonical maps
\[
  \operatorname{Ps.hom}_{S_\lambda}(X_\lambda,Y_\lambda)\to\operatorname{Ps.hom}_{S_\mu}(X_\mu,Y_\mu)
\]
\oldpage[IV-4]{240}
for $\lambda\leq\mu$, and $\operatorname{Ps.hom}_{S_\lambda}(X_\lambda,Y_\lambda)\to\operatorname{Ps.hom}_S(X,Y)$, are defined (without any separation hypothesis on the $Y_\lambda$ or on $Y$); consequently the same holds for the map \sref{IV.20.3.8.1}.
The proposition will follow from the following lemma:

\begin{lemma}[20.3.8.2]
\label{IV.20.3.8.2}
With the notation of \sref{IV.8.8.1}, suppose $X_\alpha$ quasi-compact, and let $U_\alpha$ be an open subset of $X_\alpha$; let $U_\lambda$ and $U$ be its inverse images in $X_\lambda$ and $X$ for $\lambda\geq\alpha$.
Suppose that one of the conditions~(i),~(ii) of \sref{IV.20.3.8} is satisfied.
Then, for $U$ to be schematically dense in $X$, it is necessary and sufficient that there exist $\lambda\geq\alpha$ such that $U_\lambda$ is schematically dense in $X_\lambda$, and in this case $U_\mu$ is schematically dense in $X_\mu$ for $\mu\geq\lambda$.
\end{lemma}

Suppose first that condition~(i) is satisfied; denote by $j_\lambda:U_\lambda\to X_\lambda$ and $j:U\to X$ the canonical injections, and by $\sh{J}_\lambda$ the kernel of the canonical homomorphism $\sh{O}_{X_\lambda}\to(j_\lambda)_*(\sh{O}_{U_\lambda})$.
Since the immersion $j_\lambda$ is quasi-compact and quasi-separated, $(j_\lambda)_*(\sh{O}_{U_\lambda})$ is a quasi-coherent $\sh{O}_{X_\lambda}$-module, so $\sh{J}_\lambda$ is a quasi-coherent ideal of $\sh{O}_{X_\lambda}$, and since $X_\lambda$ is Noetherian, $\sh{J}_\lambda$ is coherent, hence of finite type.
On the other hand, since the transition morphism $p_{\lambda\mu}:X_\mu\to X_\lambda$ (resp. $p_\lambda:X\to X_\lambda$) is flat, it follows from \sref{IV.2.3.1} that we can identify $(j_\mu)_*(\sh{O}_{U_\mu})$ with $p_{\lambda\mu}^*((j_\lambda)_*(\sh{O}_{U_\lambda}))$ (resp. $j_*(\sh{O}_U)$ with $p_\lambda^*((j_\lambda)_*(\sh{O}_{U_\lambda}))$).
The assertion then follows from the definition of a schematically dense open subset and from \sref{IV.8.5.8}[(ii)].

To establish \sref{IV.20.3.8.2} when condition~(ii) is satisfied, we will first prove two lemmas.

\begin{lemma}[20.3.8.3]
\label{IV.20.3.8.3}
Under the hypotheses of \sref{IV.8.2.2}, let $M_\lambda$ (resp. $M$) be the set of maximal points of $S_\lambda$ (resp. $S$).
Suppose that for every $\lambda$ the set $M_\lambda$ is finite, and that the $M_\lambda$ form a projective system of sets.
Then $M$ is the projective limit of the system of the $M_\lambda$.
\end{lemma}

Let us first show that a point $s\in\varprojlim M_\lambda$ is maximal in $S$: indeed, if $s'$ is a generization of $s$, the image $s'_\lambda$ of $s'$ in $S_\lambda$ is a generization of the image $s_\lambda$ of $s$, hence is equal to $s_\lambda$ for every $\lambda$, which implies $s'=s$, since the underlying set of $S$ is the projective limit of the underlying sets of the $S_\lambda$ \sref{IV.8.2.9}.
Conversely, let $s$ be a maximal point of $S$, and let us prove that it belongs to $\varprojlim M_\lambda$.
Let $s_\lambda$ be the image of $s$ in $S_\lambda$, and $M'_\lambda$ the set of maximal points of $S_\lambda$ that are generizations of $s_\lambda$; the $M'_\lambda$ are non-empty finite sets which form a projective system, so $M'=\varprojlim M'_\lambda$ is not empty and the map $M'\to M'_\lambda$ is surjective (Bourbaki, \emph{Ens.}, chap.~III, 2\textsuperscript{e} éd., \textsection~7, n\textsuperscript{o}~4, Example~I).
On the other hand, we have $\Spec(\sh{O}_{S,s})=\varprojlim\Spec(\sh{O}_{S_\lambda,s_\lambda})$ by virtue of \sref{IV.8.2.12} and \sref{IV.8.2.9}, so the points of $\varprojlim M'_\lambda$ are also maximal points of $\Spec(\sh{O}_{S,s})$ by the first part of the argument.
Hence $M'=\varprojlim M'_\lambda$ necessarily reduces to the point $s$; we conclude that $M'_\lambda$ reduces to the point $s_\lambda$, and consequently the $s_\lambda$ are maximal in the $S_\lambda$, which completes the proof of the lemma.

\begin{lemma}[20.3.8.4]
\label{IV.20.3.8.4}
The hypotheses being those of \sref{IV.20.3.8.3}, suppose moreover $S_\alpha$ quasi-compact; let $V_\alpha$ be an open set of $S_\alpha$, and let $V_\lambda$ and $V$ be its inverse images in $S_\lambda$ for $\lambda\geq\alpha$ and in $S$.
If $U_\alpha$ is dense in $S_\alpha$, then $U_\lambda$ is dense in $S_\lambda$ for $\lambda\geq\alpha$ and $U$ is dense in $S$.
Conversely, if $U$ is dense in $S$ and if moreover the set $M$ of maximal points of $S$ is finite, there exists $\lambda\geq\alpha$ such that $U_\lambda$ is dense in $S_\lambda$.
\end{lemma}

Indeed, since $M_\alpha$ is finite, the hypothesis that $U_\alpha$ is dense in $S_\alpha$ implies that
\oldpage[IV-4]{241}
$M_\alpha\subset U_\alpha$, hence $M_\lambda\subset U_\lambda$ for $\lambda\geq\alpha$ and $M\subset U$ by virtue of \sref{IV.20.3.8.3}, which proves the first assertion.
Conversely, suppose $M$ finite and $U$ dense in $S$, so that $M\subset U$; note that the $U_\lambda$ are open, hence ind-constructible, and the $M_\lambda$ are finite, hence pro-constructible \sref{IV.1.9.6}.
The second assertion then follows from \sref{IV.8.3.2}.

\phantomsection\egalabel{IV.20.3.8.5}{20.3.8.5}
\emph{20.3.8.5. End of the proof of \sref{IV.20.3.8.2}.}
Suppose condition~(ii) is satisfied, and note that $X$ is then \emph{reduced} by virtue of \sref{IV.8.7.1}; it therefore amounts to the same thing to say that $U_\lambda$ (resp. $U$) is schematically dense in $X_\lambda$ (resp. $X$) or that it is dense in $X_\lambda$ (resp. $X$), and the conclusion follows from \sref{IV.20.3.8.4} applied to the $X_\lambda$ and to $X$.

\phantomsection\egalabel{IV.20.3.8.6}{20.3.8.6}
\emph{20.3.8.6. End of the proof of \sref{IV.20.3.8}.}
To show that the map \sref{IV.20.3.8.1} is injective, consider two morphisms $u'_\lambda$, $u''_\lambda$ from a schematically dense open subset $U_\lambda\subset X_\lambda$ to $Y_\lambda$, and suppose that their images $u'$, $u''$, which are morphisms from $U$ to $Y$, coincide on a schematically dense open subset $V$ of $U$.
Moreover, under either of the hypotheses~(i),~(ii), we may suppose $V$ \emph{quasi-compact}; this is evident if $X$ is Noetherian; otherwise, since $X$ has only finitely many maximal points and is reduced, it suffices to replace $V$ by the union of finitely many affine open neighbourhoods of these maximal points (contained in $V$ by hypothesis).
Then there exists a $\lambda\geq\alpha$ such that $V$ is the inverse image of a quasi-compact open subset $V_\lambda$ of $X_\lambda$ \sref{IV.8.2.11}, and it follows from \sref{IV.20.3.8.2} that, by taking $\lambda$ sufficiently large, we may suppose that $V_\lambda$ is schematically dense in $X_\lambda$.
It then follows from \sref{IV.8.8.2}[(i)] that, by taking $\lambda$ sufficiently large, $u'_\lambda$ and $u''_\lambda$ coincide on $V_\lambda$, and therefore belong to the same pseudo-homomorphism.

Finally, let us prove that the map \sref{IV.20.3.8.1} is surjective.
We now consider a morphism $u$ from a schematically dense open subset $U$ of $X$ to $Y$; as above, we may suppose $U$ quasi-compact and quasi-separated (in case~(ii), $U$ can be replaced by a union of finitely many pairwise disjoint affine open subsets).
We may therefore again suppose that there exists a $\lambda\geq\alpha$ such that $U$ is the inverse image of a quasi-compact open subset $U_\lambda$ of $X_\lambda$ \sref{IV.8.2.11}, which is automatically quasi-separated, and by \sref{IV.20.3.8.2} we may moreover suppose that $U_\lambda$ is schematically dense in $X_\lambda$.
Since the $Y_\lambda$ are assumed locally of finite presentation, it follows from \sref{IV.8.8.2}[(i)] that there exists a $\lambda$ such that $u$ is the image of a morphism $u_\lambda:U_\lambda\to Y_\lambda$; whence the conclusion.

\begin{remarks}[20.3.8.7]
\label{IV.20.3.8.7}
\begin{enumerate}[label=(\roman*)]
  \item To prove that the map \sref{IV.20.3.8.1} is \emph{injective}, it is not necessary, under hypothesis~(i) of \sref{IV.20.3.8}, to suppose $X$ Noetherian.
  Indeed, Lemma \sref{IV.20.3.8.2} does not use this hypothesis.
  With the notation of \sref{IV.20.3.8.6}, let $Z_\lambda$ be the sub-prescheme of coincidences of $u'_\lambda$ and $u''_\lambda$, and let $Z$ be the sub-prescheme of coincidences of $u'$ and $u''$; it follows from definition \sref{IV.17.4.5} and from \sref[I]{I.3.3.10.1} that $Z$ is the projective limit of the $Z_\mu$ for $\mu\geq\lambda$.
  Now, by hypothesis, $Z$ dominates an open subset schematically dense in $X$; it follows that $Z$ is itself induced on an open subset of $X$, by virtue of the following lemma:

\begin{lemma}[20.3.8.8]
\label{IV.20.3.8.8}
Let $T$ be a prescheme.
Then every sub-prescheme $W$ of $T$ that dominates a schematically dense open subset $V$ of $T$ is induced on a (schematically dense) open subset of $T$.
\end{lemma}

  Indeed, the subspace of $T$ underlying $W$ is locally closed, hence open
  \oldpage[IV-4]{242}
  in its closure, which already proves that the underlying space of $W$ is open in $T$; the conclusion then follows from \sref{IV.11.10.1}[(c)].

  This lemma being established, we conclude that for $\mu\geq\lambda$ sufficiently large, $Z_\mu$ is induced on an open subset of $X_\mu$ by virtue of \sref{IV.8.6.3}, because $Z_\lambda$, being a sub-prescheme of a Noetherian prescheme, is of finite presentation over $X_\lambda$ \sref{IV.1.6.2}, and therefore so are the $Z_\mu$ over $X_\mu$ for $\mu\geq\lambda$, and $Z$ over $X$.
  We can now apply \sref{IV.20.3.8.2}, which shows that for $\mu\geq\lambda$ sufficiently large, $Z_\mu$ is schematically dense in $X_\mu$, whence the conclusion.

  \item If, in hypothesis~(i) of \sref{IV.20.3.8}, we drop the condition that $X$ is Noetherian, we see that the argument of \sref{IV.20.3.8.6} still shows that the image of \sref{IV.20.3.8.1} consists of the pseudo-$S$-morphisms having a representative which is an $S$-morphism $U\to Y$, where $U$ is schematically dense in $X$ and \emph{quasi-compact and quasi-separated}.
\end{enumerate}
\end{remarks}
\end{proof}

\begin{corollary}[20.3.9]
\label{IV.20.3.9}
Suppose that one or other of the hypotheses~a),~b) of \sref{IV.20.3.6} holds for $X$, and that $Y$ is locally of finite presentation over $S$.
Then, for a pseudo-$S$-morphism $\omega$ from $X$ to $Y$ to be defined at the point $x$ \sref{IV.20.2.3}, it is necessary and sufficient that its restriction to $\Spec(\sh{O}_{X,x})$ be everywhere defined (in other words, that it be an $S$-morphism from $\Spec(\sh{O}_{X,x})$ to $Y$).
\end{corollary}

The following result, which we will use in the proof of \sref{IV.20.3.11}, uses the theory of faithfully flat descent of chap.~VI.
The reader may check that the results of (20.3) will not be used in that theory.

\begin{lemma}[20.3.10]
\label{IV.20.3.10}
Let $f:X'\to X$ be a faithfully flat and quasi-compact $S$-morphism, $X''=X'\times_X X'$, $p_1$ and $p_2$ the canonical projections of $X''$ to $X'$, and $Y$ a prescheme separated over $S$.
Let $U$ be an open subset of $X$, $U'=f^{-1}(U)$, $U''=p_1^{-1}(U')=p_2^{-1}(U')$, and suppose that $U''$ is schematically dense in $X''$.
Let $g:U\to Y$ be an $S$-morphism; then, if $g\circ(f|U')$ extends to an $S$-morphism $X'\to Y$, $g$ extends to an $S$-morphism $X\to Y$.
\end{lemma}

Note that the hypotheses imply that $U$ (resp. $U'$) is schematically dense in $X$ (resp. $X'$) \sref{IV.11.10.5}[(i)]; we can therefore also say that, if $\omega$ denotes the pseudo-$S$-morphism that is the class of $g$, the statement of \sref{IV.20.3.10} means that \emph{if $\omega\circ f$ is everywhere defined, so is $\omega$}.

\begin{proof}
To prove \sref{IV.20.3.10}, denote by $g'$ an $S$-morphism $X'\to Y$ that extends $g\circ(f|U')$, and set $g''_i=g'\circ p_i:X''\to Y$ ($i=1,2$).
If we set $f''=f\circ p_1=f\circ p_2:X''\to X$, it is clear that $g''_1$ and $g''_2$ coincide on $U''$ with $g\circ(f''|U'')$.
But since $Y$ is separated over $S$ and $U''$ is schematically dense in $X''$, we have $g''_1=g''_2$ \sref{IV.11.10.1}[(d)].
Since $f$ is faithfully flat and quasi-compact, it follows from descent theory (chap.~VI) that there exists a unique $S$-morphism $h:X\to Y$ such that $h\circ f=g'$; since the restriction $U'\to U$ of $f$ is a faithfully flat and quasi-compact morphism and $U''=U'\times_U U'$, the preceding uniqueness result, applied to $U$ instead of $X$, shows that $h|U=g$, which proves the lemma.
\end{proof}

\begin{proposition}[20.3.11]
\label{IV.20.3.11}
Let $Y$ be an $S$-prescheme separated over $S$, $\omega$ a pseudo-$S$-morphism from $X$ to $Y$, and $f:X'\to X$ an $S$-morphism.
Suppose that $f$ is flat, and that one of the following conditions is fulfilled:
\begin{enumerate}[label=\emph{(\roman*)}]
\oldpage[IV-4]{243}
  \item $f$ is an open morphism, or is surjective and quasi-compact, and $\operatorname{dom}_S(\omega)$ contains an open subset $U$ that is schematically dense in $X$ and retrocompact in $X$.
  \item $f$ is locally of finite presentation.
  \item $Y$ is locally of finite presentation over $S$, and $X$ satisfies one of the conditions~a),~b) of \sref{IV.20.3.6}.
\end{enumerate}
Then $f^{-1}(\operatorname{dom}_S(\omega))$ is schematically dense in $X'$, so that $\omega\circ f$ is defined \sref{IV.20.3.2}, and we have
\[
\label{IV.20.3.11.1}
  \operatorname{dom}_S(\omega\circ f)=f^{-1}(\operatorname{dom}_S(\omega)).
  \tag{20.3.11.1}
\]
\end{proposition}

\begin{proof}
Let us first prove that $f^{-1}(\operatorname{dom}_S(\omega))$ is schematically dense in $X'$.
The question being local on $X$ and $X'$, we may suppose $X$ and $X'$ affine, and since $f$ is flat, it suffices to see, by virtue of \sref{IV.11.10.5}[(ii), (a)], that $\operatorname{dom}_S(\omega)$ contains an open set $U$ that is retrocompact and schematically dense in $X$.
This follows from the hypothesis in case~(i), and from \sref{IV.20.2.12} in case~(iii), taking into account the fact that in an affine scheme every open subset of the form $\mathrm{D}(t)$ is retrocompact; finally, in case~(ii), we see directly that $f^{-1}(\operatorname{dom}_S(\omega))$ is schematically dense in $X'$ by applying \sref{IV.20.3.5}[(i)].

Let us now prove \sref{IV.20.3.11.1}, in other words that, for every point $x'\in\operatorname{dom}_S(\omega\circ f)$, we have $x=f(x')\in\operatorname{dom}_S(\omega)$.
Note first that we may restrict to the case where $f$ is \emph{faithfully flat and quasi-compact}.
Indeed, this is the hypothesis in the second case of~(i); in the other cases, the question is local on $X'$, and we may therefore already suppose $X$ and $X'$ affine, hence $f$ quasi-compact.
In the first case of~(i) and in case~(ii), $f$ is an open morphism \sref{IV.11.3.1}, so, replacing $X$ by the open subset $f(X')$, we may suppose $f$ surjective, hence faithfully flat.
In case~(iii), using \sref{IV.20.3.9}, we may restrict to proving that the restriction of $\omega$ to $\Spec(\sh{O}_{X,x})$ is everywhere defined, and we can therefore replace $X$ by $\Spec(\sh{O}_{X,x})$, $X'$ by $X'\times_X\Spec(\sh{O}_{X,x})$ and $f$ by its restriction to this last prescheme, which is a surjective morphism \sref{IV.2.3.4}, hence faithfully flat.

So suppose $f$ faithfully flat and quasi-compact; with the notation of Lemma \sref{IV.20.3.10}, it then suffices to see that $U''$ is schematically dense in $X''$, taking for $U$ an open subset schematically dense in $X$ and contained in $\operatorname{dom}_S(\omega)$; this will be the case, by virtue of \sref{IV.11.10.5}[(ii), (a)], if $U$ is taken retrocompact in $X$ (since the morphism $f'':X''\to X$ is flat).
Now the existence of such an open subset $U$ is part of the hypothesis in case~(i); in case~(iii) it follows from \sref{IV.20.2.12} and from the fact that in an affine scheme $\Spec(A)$ every open subset of the form $\mathrm{D}(t)$ is retrocompact.
Finally, in case~(ii), take $U=\operatorname{dom}_S(\omega)$ and let us show directly that $U''$ is schematically dense in $X''$: for this it suffices to remark that $f'':X''\to X$ is flat and locally of finite presentation and to apply \sref{IV.11.10.5}[(ii), (b)].
\end{proof}

\begin{corollary}[20.3.12]
\label{IV.20.3.12}
Let $\vphi$ be a pseudo-function on a prescheme $X$.
Then, for every flat morphism $f:X'\to X$ locally of finite presentation, the pseudo-function $\vphi\circ f$ is defined and we have $\operatorname{dom}(\vphi\circ f)=f^{-1}(\operatorname{dom}(\vphi))$.
\end{corollary}

\begin{remark}[20.3.13]
\label{IV.20.3.13}
\oldpage[IV-4]{244}
When $X$ satisfies one of the conditions~a),~b) of \sref{IV.20.3.6}, we have seen \sref{IV.20.2.11} that the pseudo-functions on $X$ are identified with the meromorphic functions on $X$.
By virtue of \sref{IV.20.1.12} and \sref{IV.20.2.13}[(i)], if we only suppose that the morphism $f:X'\to X$ is \emph{flat}, then, for every pseudo-function $\vphi$ on $X$, $\vphi\circ f$ is defined and we have $\operatorname{dom}(\vphi\circ f)=f^{-1}(\operatorname{dom}(\vphi))$.
\end{remark}


\subsection{Properties of the domains of definition of rational functions}
\EGAReaderTarget{ega4.c04.subsec.20-4}{IV.20.4}

\begin{env}[20.4.1]\label{IV.20.4.1}
\EGAReaderTarget{ega4.c04.statement.20-4-1}{IV.20.4.1}
Let $X$, $Y$ be two $S$-preschemes, $\omega$ a pseudo-$S$-morphism from $X$ to $Y$.
Let $u$ be an $S$-morphism $U\to Y$ belonging to $\omega$, where $U$ is schematically dense in $X$, and consider the \emph{graph} $\Gamma_u$, a sub-prescheme of $U\times_S Y$ \sref[I]{I.5.3.11}, hence a sub-prescheme of $X\times_S Y$.
Suppose that this sub-prescheme has a \emph{closure} \sref[I]{I.9.5.11} in $X\times_S Y$; if $j:\Gamma_u\to X\times_S Y$ is the canonical injection, this will be the case when the $\sh{O}_{X\times_S Y}$-module $j_*(\sh{O}_{\Gamma_u})$ is quasi-coherent; it follows from the definition of the equivalence class $\omega$ \sref{IV.20.2.1} that $j_*(\sh{O}_{\Gamma_u})$ does not depend on the representative $u$ considered, so the closure prescheme of $\Gamma_u$ is then well determined by $\omega$; we say that this closed sub-prescheme of $X\times_S Y$ is the \emph{graph} of the pseudo-$S$-morphism $\omega$, and we denote it by $\Gamma_\omega$.
Note that $\Gamma_\omega$ is defined if there exists a morphism $u:U\to Y$ of the class $\omega$ such that $U$ is \emph{retrocompact} in $X$ and if moreover $Y$ is quasi-separated over $S$, because the injection $j$ is then a quasi-compact and quasi-separated morphism (\sref[I]{I.2.2} and \sref[I]{subsection:I.7.4}); the first hypothesis will always be satisfied when $X$ is \emph{locally Noetherian}.
Note also that when $X$ is \emph{reduced}, so is $\Gamma_u$, which is isomorphic to $U$ \sref[I]{I.5.3.11}; then $\Gamma_\omega$ exists and is none other than the \emph{reduced} sub-prescheme of $X\times_S Y$ having as underlying space the closure in $X\times_S Y$ of the space underlying $\Gamma_u$ (\sref[I]{I.5.2.1} and \sref[I]{I.5.2.2}).
Note finally that if $Y$ is \emph{separated} over $S$, $\Gamma_u$ is \emph{closed} in $U\times_S Y$ \sref[I]{I.5.4.3}, hence is \emph{induced by} $\Gamma_\omega$ (when the latter exists) on the open subset $\Gamma_\omega\cap(U\times_S Y)$ \sref[I]{I.9.5.10}; on the other hand, this induced prescheme is in general different from $\Gamma_u$ when $Y$ is not separated.
In particular, if $v:X\to Y$ is an $S$-morphism, the graph $\Gamma_\omega$ of the class $\omega$ of $v$ may be distinct from the graph $\Gamma_v$.
Therefore, in what follows, whenever the graph of a pseudo-$S$-morphism is in question, we will restrict ourselves to the case where \emph{$Y$ is separated over $S$}.
\end{env}

\begin{env}[20.4.2]
\label{IV.20.4.2}
So suppose $\Gamma_\omega$ defined and $Y$ separated over $S$; denote by $p$ and $q$ the restrictions to $\Gamma_\omega$ of the canonical projections
\[
  \xymatrix{
  & \Gamma_\omega \ar[dl]_p \ar[dr]^q & \\
  X & & Y
  }
\]
Then, if $U\subset\operatorname{dom}_S(\omega)$, the restriction $p^{-1}(U)\to U$ of $p$ is an isomorphism \sref[I]{I.5.3.11}; conversely, if $U$ is an open subset of $X$ having this property, and if $u$ is the isomorphism inverse to the restriction $p^{-1}(U)\to U$ of $p$, then $q\circ u$ is an $S$-morphism from $U$ to $Y$ which coincides with a morphism of the class $\omega$ on $U\cap\operatorname{dom}_S(\omega)$.
We conclude that $\operatorname{dom}_S(\omega)$ is the \emph{largest open subset} $U$ of $X$ such that the restriction $p^{-1}(U)\to U$ of $p$ is an
\oldpage[IV-4]{245}
isomorphism.
Let $v:\operatorname{dom}_S(\omega)\to\Gamma_\omega$ be the $S$-morphism inverse to the restriction $p^{-1}(\operatorname{dom}_S(\omega))\to\operatorname{dom}_S(\omega)$ of $p$; we sometimes denote by $p^{-1}$ the pseudo-$S$-morphism from $X$ to $\Gamma_\omega$ that is the class of $v$ (whose associated rational map \sref{IV.20.2.13}[(ii)] is then \emph{birational}); since $p^{-1}(\operatorname{dom}_S(\omega))$ is the graph of the $S$-morphism $u=q\circ v:\operatorname{dom}_S(\omega)\to Y$, it is schematically dense in $\Gamma_\omega$ \sref{IV.11.10.3}[(iv)], so $\omega$ can be regarded as the \emph{composite} $q\circ p^{-1}$ \sref{IV.20.3.2}.
For every subset $M$ of the space underlying $X$, we sometimes set $\omega(M)=q(p^{-1}(M))$, which then amounts to regarding $\omega$ as a map from $X$ to $\mathfrak{P}(Y)$ (or, as some authors say, a ``multivalued function'').
Note that when $x\in\operatorname{dom}_S(\omega)$, $\omega(\{x\})$ is the set $\{u(x)\}$; in general, for an arbitrary $x\in X$, if we denote by $s$ the image of $x$ in $S$ and by $Y_s$ the fibre at the point $s$ of the structural morphism $Y\to S$, the set $\omega(\{x\})$ is a subset of the prescheme $Y_s$.
\end{env}

\begin{env}[20.4.3]
\label{IV.20.4.3}
In all the rest of this section, we will restrict ourselves to the case where $X$ is \emph{reduced}, so that pseudo-$S$-morphisms are then the same as \emph{rational $S$-maps} \sref{IV.20.2.7}.
Moreover, the \emph{graph} of every rational $S$-map from $X$ to $Y$ is then defined \sref{IV.20.4.1}.
\end{env}

\begin{proposition}[20.4.4]
\label{IV.20.4.4}
Let $X$ be a locally integral $S$-prescheme, $Y$ an $S$-prescheme locally of finite type and separated over $S$, $\omega$ a rational $S$-map from $X$ to $Y$, and $p:\Gamma_\omega\to X$ the canonical projection.
Then, if $x\in X$ is a normal point such that the set $p^{-1}(x)$ contains an isolated point, $\omega$ is defined at the point $x$.
\end{proposition}

\begin{proof}
Indeed, the first projection $p_1:X\times_S Y\to X$ is then a separated morphism locally of finite type, hence so is its restriction $p:\Gamma_\omega\to X$, which is moreover \emph{birational}; and since $\Gamma_\omega$ is reduced and $X$ integral, $\Gamma_\omega$ is \emph{integral}; it therefore follows from (\textbf{Err}\textsubscript{\textbf{IV}}, 30) that the hypothesis on $x$ implies the existence of an open neighbourhood $U$ of $x$ such that the restriction $p^{-1}(U)\to U$ of $p$ is an isomorphism; whence the conclusion \sref{IV.20.4.2}.

Note that the statement \sref{IV.20.4.4} is the original formulation that Zariski gave of his ``Main theorem'' (except that he restricted himself to algebraic schemes over a base field).
\end{proof}

\begin{proposition}[20.4.5]
\label{IV.20.4.5}
The hypotheses and notation being those of \sref{IV.20.4.4}, suppose moreover $X$ normal, and let $X'$ be a reduced prescheme and $f:X'\to X$ a morphism locally of finite type and universally open.
Then $\omega\circ f$ is defined, and we have $\operatorname{dom}_S(\omega\circ f)=f^{-1}(\operatorname{dom}_S(\omega))$ (in other words, if $x'\in X'$ and $x=f(x')$, then, for $\omega$ to be defined at the point $x$, it is necessary and sufficient that $\omega\circ f$ be defined at the point $x'$).
\end{proposition}

\begin{proof}
Since $X'$ is reduced and $f^{-1}(\operatorname{dom}_S(\omega))$ is everywhere dense in $X'$ by virtue of \sref{IV.2.4.11}, the composite $\omega\circ f$ is defined; to prove that $\omega$ is defined at the point $x$ when $\omega\circ f$ is defined at the point $x'$, we may evidently replace $X'$ by an open neighbourhood of $x'$, and hence suppose $\omega\circ f$ \emph{everywhere defined}; moreover, since $f(X')$ is open in $X$, we may suppose $f$ \emph{surjective}.
By virtue of the hypothesis on $f$, the morphism $f_{(Y)}:X'\times_S Y\to X\times_S Y$ is then open, so $\overline{f_{(Y)}^{-1}(M)}=f_{(Y)}^{-1}(\overline{M})$ for every subset $M$ of $X\times_S Y$; taking for $M$ the set $\Gamma_u$, where $u:\operatorname{dom}_S(\omega)\to Y$ is the restriction of $\omega$ to $\operatorname{dom}_S(\omega)$, it follows from the preceding relation and from \sref[I]{I.3.3.12} that the set underlying $\Gamma_{\omega\circ f}$ is equal to $f_{(Y)}^{-1}(\Gamma_\omega)$; since we know
\oldpage[IV-4]{246}
that $\Gamma_{\omega\circ f}$ is a reduced prescheme \sref{IV.20.4.1}, we see that the $X'$-prescheme $\Gamma_{\omega\circ f}$ is equal to $(\Gamma_\omega\times_X X')_\mathrm{red}$.
But since by hypothesis the composite morphism $\Gamma_{\omega\circ f}\to\Gamma_\omega\times_X X'\xrightarrow{p^{(X')}}X'$ is an isomorphism, $p^{(X')}$ is necessarily \emph{radicial}; since $f$ is surjective, the same holds for $p:\Gamma_\omega\to X$ \sref[I]{I.3.4.8}, so $p^{-1}(x)$ is a set consisting of a single point \sref[I]{I.3.6.4}; it then follows from \sref{IV.20.4.4} that $\omega$ is defined at the point $x$.
\end{proof}

The following proposition gives a \emph{valuative criterion} for a rational map to be defined at a point:

\begin{proposition}[20.4.6]
\label{IV.20.4.6}
Let $S$ be a prescheme, $X$, $Y$ two $S$-preschemes; suppose $X$ locally Noetherian and $Y$ locally of finite type over $S$.
Let $U$ be an open subset dense in $X$, $h:U\to Y$ an $S$-morphism, $x\in X-U$ a normal point of $X$, and $h'_x:\Spec(k(x))\to Y$ an $S$-morphism.
In order that $h$ can be extended to an $S$-morphism $h':U'\to Y$, where $U'$ is an open subset of $X$ containing $U$ and $x$, and where the composite morphism $\Spec(k(x))\to U'\xrightarrow{h'}Y$ is the given $S$-morphism $h'_x$, it is necessary and sufficient that the following condition be satisfied:
\begin{quote}
\emph{(P) For every spectrum $S_1$ of a discrete valuation ring, with closed point $a$ and generic point $b$, and every morphism $u:S_1\to X$ such that $u(a)=x$, $u(b)\in U$, the composite morphism $h_1=h\circ(u|\{b\}):\{b\}=u^{-1}(U)\to Y$ extends to a morphism $h'_1:S_1\to Y$ such that the diagram}
\[
  \xymatrix{
  \Spec(k(a)) \ar[r] & S_1 \ar[d]^{h'_1} \\
  \Spec(k(x)) \ar[r]_{h'_x} & Y
  }
\]
\emph{is commutative.}
\end{quote}
Moreover, if this condition is satisfied, and if $h'':U''\to Y$ is a morphism satisfying the same conditions as $h'$, then there exists an open set containing $U\cup\{x\}$ on which $h'$ and $h''$ coincide.
\end{proposition}

\begin{proof}
Let us first prove the last assertion; we may suppose $h'$ and $h''$ defined on the same open subset $U_0\supset U\cup\{x\}$.
The sub-prescheme $Z$ of coincidences of $h'$ and $h''$ \sref{IV.17.4.5} contains $U$ and $x$, so there is an open neighbourhood $V$ of $x$ in $U_0$ such that the sub-prescheme induced by $Z$ on the open subset $Z\cap V$ is a \emph{closed} sub-prescheme of the prescheme induced by $X$ on $V$; since this prescheme $Z\cap V$ dominates the sub-prescheme induced by $V$ on the open subset $U\cap V$, and the latter is schematically dense in $V$, $Z\cap V$ is necessarily equal to $V$ \sref{IV.20.3.8.8}, which proves that $h'$ and $h''$ coincide on $U\cup V$.

The necessity of the condition of the statement being evident, let us prove that it is sufficient.
By virtue of the bijective correspondence between $S$-morphisms from $X$ to $Y$ and $X$-sections of $X\times_S Y$ \sref[I]{I.3.3.15}, we may restrict to the case where $S=X$, so that $h$ is a $U$-section of $Y$; note that $Y$ is then locally Noetherian, and we may evidently (since $X$ is locally integral and locally Noetherian) suppose $X=S$ irreducible.
Moreover, taking into account \sref{IV.20.3.7}, we may suppose that $X=\Spec(A)$, where $A$ is a Noetherian \emph{local} ring that is integral and integrally closed.

\oldpage[IV-4]{247}
Note that for every $x'\in U$, $h(x')$ is a specialization of $h'_x(x)=y$.
Indeed, there exist a spectrum $S_1$ of a discrete valuation ring and a morphism $u:S_1\to X$ such that $u(a)=x$, $u(b)=x'$ \sref[II]{II.7.1.9}.
Applying the condition of the statement, we immediately obtain our assertion, since $h(x')=h_1(b)=h'_1(b)$ and $y=h'_1(a)$.
If $Y'$ is an affine open neighbourhood of $y$, we therefore have $h(X\cap U)\subset Y'$; consequently we can replace $Y$ by $Y'$, in other words suppose $Y$ \emph{affine}, hence separated over $X$.
Let $\omega$ be the rational $X$-section of $Y$ to which $h$ belongs, so that its graph here has as underlying set the closure of $h(U)$ in $Y$.
Since $Y$ is separated over $X$, we can apply \sref{IV.20.4.4} to $\omega$: it will suffice to prove that, if $p:\Gamma_\omega\to X$ is the canonical projection, $p^{-1}(x)$ consists of a single point $y$ and that $h'_x(x)=y$.
Indeed, by \sref{IV.20.4.4}, $h$ will then extend to a section $h'$ over an open subset $U'$ containing $U\cup\{x\}$, such that $h'(x)=y$, and since there is then only one $X$-morphism $\Spec(k(x))\to Y$ sending $x$ to $y$, this will prove that $h'_x$ is identical with the composite of $h'$ and $\Spec(k(x))\to U'$.

Since $h(x')$ is a specialization of $y$ for $x'\in X\cap U$, we have $y\in p^{-1}(x)$.
Suppose then that $z$ is an arbitrary point of $p^{-1}(x)$.
Since $\Gamma_\omega$ is the closure of $h(U)$ and $h(U)$ consists of points lying in the closure of $h(\xi)$, where $\xi$ is the generic point of $X$, $\Gamma_\omega$ is the closure of $h(\xi)$ in $Y$.
We then take a spectrum $S_1$ of a discrete valuation ring and a morphism $v:S_1\to Y$ such that $v(a)=z$, $v(b)=h(\xi)$ \sref[II]{II.7.1.9}, and set $u=p\circ v$, so that $u(a)=x$, $u(b)=\xi$.
Applying the condition of the statement to $u$, we see that we obtain a morphism $h'_1:S_1\to Y$ such that $h'_1(a)=y$ and $h'_1(b)=h(\xi)$; but this implies $z=y$ by virtue of \sref[II]{II.7.2.3}, since $Y$ is separated over $X$, so that $v$ and $h'_1$ must coincide, because they are equal at the point $b$.
Q.E.D.
\end{proof}

\begin{corollary}[20.4.7]\label{IV.20.4.7}
\EGAReaderTarget{ega4.c04.cor.20-4-7}{IV.20.4.7}
The hypotheses on $S$, $X$, $Y$, $U$ and $h$ being those of \sref{IV.20.4.6}, let $E$ be a subset of $X-U$ such that $X$ is normal at every point of $E$, and for each $x\in E$ let $h'_x:\Spec(k(x))\to Y$ be an $S$-morphism such that condition~(P) of \sref{IV.20.4.6} is satisfied.
Suppose moreover that the union $F$ of the graphs of the $h'_x$ (identified with subsets of $X\times_S Y$) for $x\in E$ is contained in the union of finitely many open subsets $V_i$ of $X\times_S Y$ which are \emph{separated} over $X$ (a condition automatically satisfied if $Y$ is \emph{separated} over $S$, or if $X$ is \emph{quasi-compact} and $Y$ of \emph{finite type} over $S$).
Then there exists an $S$-morphism $h':U'\to Y$, where $U'$ is an open subset of $X$ containing $U\cup E$, such that, for every $x\in E$, the composite
\[
  \Spec(k(x))\to U'\xrightarrow{h'}Y
\]
is equal to $h'_x$.
\end{corollary}

\begin{proof}
Note first that, in \sref{IV.20.4.6}, if $Y$ is assumed \emph{separated}, there is a largest open subset $U_0\supset U$, equal to the domain of the rational $S$-map corresponding to $h$, to which $h$ can be extended, and this extension is unique \sref[I]{I.7.2.2}; whence the conclusion in this case, thanks to \sref{IV.20.4.6}.
In the general case, let $E_i$ be the set of $x\in E$ such that $(x,h'_x(x))\in V_i$.
By virtue of \sref{IV.20.4.6}, for every $x\in E$ there is an extension $h^{(x)}$ of $h$ to $U\cup W^{(x)}$, where $W^{(x)}$ is a neighbourhood of $x$ in $X$ such that the graph of $h^{(x)}|W^{(x)}$ is contained in \emph{all} the $V_i$ such that $x\in E_i$.
Since $V_i$ is separated over $X$ and $X$ is reduced, for two points $x'$, $x''$ of $E_i$, $h^{(x')}$ and $h^{(x'')}$ coincide on $W^{(x')}\cap W^{(x'')}$, since they
\oldpage[IV-4]{248}
coincide on the everywhere dense open subset $U\cap W^{(x')}\cap W^{(x'')}$ of this set.
There is therefore a morphism $h_i:U\cup U_i\to Y$ which extends $h$ to an open subset $U\cup U_i$ containing $E_i$; moreover, for every pair of indices $i$, $j$, the graphs of the restrictions $h_i|(U_i\cap U_j)$ and $h_j|(U_i\cap U_j)$ are contained in $V_i\cap V_j$; since $V_i\cap V_j$ is separated over $X$ and the preceding morphisms coincide on an everywhere dense open subset $U\cap U_i\cap U_j$ of $U_i\cap U_j$, they are equal.
The morphism $h'$ equal to $h$ on $U$ and to $h_i$ on each of the $U_i$ has the required properties.
\end{proof}

\begin{remark}[20.4.8]
\label{IV.20.4.8}
When $E=X-U$, it is clear that if, for every affine open subset $T$ of $X$, there exists an $S$-morphism $h'_T:T\to Y$ extending $h|(U\cap T)$ and such that the composite $\Spec(k(x))\to T\xrightarrow{h'_T}Y$ is equal to $h'_x$ for every $x\in T-(U\cap T)$, then the $h'_T$ are the restrictions of an (everywhere defined) $S$-morphism $h':X\to Y$, by virtue of the uniqueness assertion in \sref{IV.20.4.6}.
To prove the existence of $h'$, we are therefore reduced to the case where $X$ is \emph{affine}, and then, in order to be able to apply \sref{IV.20.4.7}, it suffices that the set $F$, the union of the graphs of the $h'_x$, be \emph{quasi-compact} in $X\times_S Y$.
This will be so when the $h'_x$ are of the form $\Spec(k(x))\to Z\xrightarrow{h_0}Y$, where $Z$ is a closed sub-prescheme of $X$ having $X-U$ as underlying space, and $h_0$ is an $S$-morphism.
\end{remark}

\begin{corollary}[20.4.9]
\label{IV.20.4.9}
Let $S$ be a locally Noetherian prescheme, $X$ a locally Noetherian prescheme, $f:X\to S$ a flat morphism, and $g:Y\to S$ a morphism locally of finite type.
Let $U$ be an open subset dense in $S$, $h:f^{-1}(U)\to Y$ an $S$-morphism, $T=S-U$, $Z$ the reduced closed sub-prescheme of $X$ having $f^{-1}(T)$ as underlying space, and $h_0:Z\to Y$ an $S$-morphism.
Suppose $X$ normal at all the points of $Z$.
In order that there exist an $S$-morphism (necessarily unique) $h':X\to Y$ extending $h$ and $h_0$, it is necessary and sufficient that the following condition be satisfied:

\emph{For every spectrum $S_1$ of a discrete valuation ring, with closed point $a$ and generic point $b$, and every morphism $u:S_1\to S$ such that $u(a)\in T$ and $u(b)\in U$, there exists an $S_1$-morphism $h'_1:X_{(S_1)}\to Y_{(S_1)}$ which extends $h_{(S_1)}:f_{(S_1)}^{-1}(b)\to Y_{(S_1)}$ and is such that the diagram}
\[
  \xymatrix{
  \Spec(k(z_1)) \ar[r] \ar[d] & Z_{(S_1)} \ar[d]^{(h_0)_{(S_1)}} \\
  X_{(S_1)} \ar[r]_{h'_1} & Y_{(S_1)}
  }
\]
\emph{is commutative for every $z_1\in Z_{(S_1)}$.}

Indeed, the hypothesis that $f$ is flat implies that $f^{-1}(U)$ is dense in $X$ \sref{IV.2.3.10}, and it then suffices to apply \sref{IV.20.4.8}.
\end{corollary}

\begin{corollary}[20.4.10]
\label{IV.20.4.10}
Under the hypotheses of \sref{IV.20.4.6}, suppose moreover $Y$ separated and locally quasi-finite over $S$.
Let $U$ be an open subset dense in $X$, $h:U\to Y$ an $S$-morphism, $\omega$ the corresponding rational $S$-map, and $x\in X-U$ a normal point of $X$.
In order that $\omega$ be defined at the point $x$, it is necessary and sufficient that the following condition be satisfied: there exist a spectrum $S_1$ of a discrete valuation ring, with closed point $a$ and generic point $b$, and a morphism $u:S_1\to X$
\oldpage[IV-4]{249}
such that $u(a)=x$, $u(b)\in U$ and the rational $S$-map $\omega\circ u$ is everywhere defined.

Indeed, by hypothesis all the fibres of the projection morphism $X\times_S Y\to X$ consist of isolated points, and to apply \sref{IV.20.4.4} it suffices to know that the fibre $p^{-1}(x)$ in $\Gamma_\omega$ is non-empty.
Now if $h_1$ is the unique morphism of the class $\omega\circ u$, the unique point of $X\times_S Y$ above $x$ and $h_1(a)$ belongs to $\Gamma_\omega$, whence the conclusion.
\end{corollary}

\begin{proposition}[20.4.11]
\label{IV.20.4.11}
Let $X$ be a locally Noetherian prescheme, $Y$ an $S$-prescheme affine over $S$, $U$ an open subset of $X$, and $Z=X-U$.
Suppose that $\operatorname{prof}_Z(\sh{O}_X)\geq 2$ \sref{IV.5.10.1}; then every $S$-morphism $f:U\to Y$ extends uniquely to an $S$-morphism from $X$ to $Y$.
\end{proposition}

\begin{proof}
We may restrict to the case where $S$ and $X$ are affine and (by virtue of \sref[I]{I.3.3.14}) to the case where $S=X$; we then have $X=\Spec(A)$, $Y=\Spec(B)$, $B$ being an $A$-algebra of finite type.
Since $B$ is a quotient of a polynomial algebra $A[(T_\lambda)]_{\lambda\in L}$, $Y$ is a closed sub-prescheme of $Y'=X[(T_\lambda)]_{\lambda\in L}$.
On the other hand, the hypothesis on $Z$ implies that $U$ is schematically dense in $X$, by virtue of \sref{IV.20.2.13}[(iv)] and \sref{IV.5.10.2}.
If we prove that every $X$-morphism $u:U\to Y$ extends uniquely to an $X$-morphism $v':X\to Y'$, it will follow that $v'$ factors as $X\xrightarrow{v}Y\to Y'$; indeed, the sub-prescheme $v'^{-1}(Y)$ is closed and dominates $U$ \sref[I]{I.4.4.1}, hence is identical with $X$ \sref{IV.20.3.8.8}.
Under these conditions, $v$ will be the unique $S$-morphism from $X$ to $Y$ extending $u$.
We may therefore restrict to the case where $Y=Y'$.
But then there is a bijective correspondence between the $X$-morphisms from an open subset $U\subset X$ to $Y'$ and the families $(s_\lambda)_{\lambda\in L}$ of sections of $\sh{O}_X$ over $U$ \sref[II]{II.1.7.9}; the conclusion therefore follows from \sref{IV.5.10.5}.
\end{proof}

\begin{corollary}[20.4.12]
\label{IV.20.4.12}
Let $X$ be a locally Noetherian reduced $S$-prescheme satisfying condition $(S_2)$ \sref{IV.5.7.2} (for example a locally Noetherian normal $S$-prescheme \sref{IV.5.8.6}), $Y$ an $S$-prescheme affine over $S$, and $f$ a rational $S$-map from $X$ to $Y$; then every irreducible component of $X-\operatorname{dom}(f)$ is of codimension $1$ in $X$.
\end{corollary}

\begin{proof}
It amounts to the same thing to say that if $Z^{(2)}$ is the set of $x\in X$ such that $\dim(\sh{O}_{X,x})\geq 2$, then, for every closed subset $Z\subset Z^{(2)}$ of $X$, every $S$-morphism from $X-Z$ to $Y$ extends to an $S$-morphism from $X$ to $Y$; now this follows from the hypothesis on $X$ \sref{IV.5.7.2} and from \sref{IV.20.4.11}.
\end{proof}

\subsection{Relative pseudo-morphisms}
\EGAReaderTarget{ega4.c04.subsec.20-5}{IV.20.5}

\begin{env}[20.5.1]\label{IV.20.5.1}
\EGAReaderTarget{ega4.c04.statement.20-5-1}{IV.20.5.1}
Let $X$, $Y$ be two $S$-preschemes.
It follows from the definitions \sref{IV.11.10.8} that the intersection of two open subsets $U$, $U'$ of $X$ that are \emph{universally schematically dense with respect to $S$} again has this property.
We therefore define an equivalence relation between $S$-morphisms $u:U\to Y$ by replacing, in \sref{IV.20.2.1}, ``schematically dense'' by ``universally schematically dense with respect to $S$''.
An equivalence class for this relation is called a \emph{pseudo-morphism from $X$ to $Y$ relative to $S$}, and the set of these classes is denoted by $\operatorname{Ps.hom}_{X/S}(X,Y)$.
\end{env}

\begin{env}[20.5.2]
\label{IV.20.5.2}
Since every open subset of $X$ universally schematically dense with respect to $S$ is in particular schematically dense, the elements of a pseudo-morphism
\oldpage[IV-4]{250}
from $X$ to $Y$ relative to $S$ are equivalent in the sense of \sref{IV.20.2.1}, whence a canonical map
\[
\label{IV.20.5.2.1}
  \operatorname{Ps.hom}_{X/S}(X,Y)\to\operatorname{Ps.hom}_S(X,Y).
  \tag{20.5.2.1}
\]
\end{env}

\begin{proposition}[20.5.3]
\label{IV.20.5.3}
Suppose $Y$ separated over $S$.
Then:
\begin{enumerate}[label=\emph{(\roman*)}]
  \item The map \sref{IV.20.5.2.1} is injective and identifies $\operatorname{Ps.hom}_{X/S}(X,Y)$ with the subset of $\operatorname{Ps.hom}_S(X,Y)$ consisting of the pseudo-$S$-morphisms $\omega$ such that $\operatorname{dom}_S(\omega)$ is universally schematically dense relative to $S$.
  \item The presheaf $U\mapsto\operatorname{Ps.hom}_{U/S}(U,Y)$ on $X$ is a sheaf.
\end{enumerate}
\end{proposition}

\begin{proof}
Assertion~(i) is immediate, because saying that two $S$-morphisms $u:U\to Y$, $u':U'\to Y$ are equivalent in the sense of \sref{IV.20.2.1} then means that they are the restrictions of one and the same morphism $\operatorname{dom}_S(\omega)\to Y$ \sref{IV.20.2.4}, and if $U$ and $U'$ are universally schematically dense with respect to $S$, then \emph{a fortiori} so is $\operatorname{dom}_S(\omega)$.
To prove~(ii), note that $U\mapsto\operatorname{Ps.hom}_S(U,Y)$ is then a sheaf \sref{IV.20.2.6}; on the other hand, given an open covering $(U_\alpha)$ of $U$ and a pseudo-$S$-morphism $\omega$ from $U$ to $Y$, it follows immediately from the definitions (cf.\ \sref{IV.20.2.1}) that, for $\operatorname{dom}_S(\omega)$ to be universally schematically dense in $U$ relative to $S$, it is necessary and sufficient that each of the sets $\operatorname{dom}_S(\omega)\cap U_\alpha=\operatorname{dom}_S(\omega|U_\alpha)$ be universally schematically dense in $U_\alpha$ relative to $S$; whence~(ii).
\end{proof}

When $Y$ is separated, we will denote by $\mathscr{P}\!s.\!\operatorname{hom}_{X/S}(X,Y)$ the sub-sheaf
\[
  U\mapsto\operatorname{Ps.hom}_{U/S}(U,Y)
\]
of $\mathscr{P}\!s.\!\operatorname{hom}_S(X,Y)$.

When $X$ is \emph{flat and locally of finite presentation} over $S$ and $Y$ is separated over $S$, the pseudo-morphisms from $X$ to $Y$ relative to $S$ can also be identified with the pseudo-$S$-morphisms $\omega$ such that, for every fibre $X_s$ of the morphism $X\to S$, $\operatorname{dom}_S(\omega)\cap X_s$ is \emph{schematically dense} in $X_s$ \sref{IV.11.10.10}.

\begin{env}[20.5.4]\label{IV.20.5.4}
\EGAReaderTarget{ega4.c04.statement.20-5-4}{IV.20.5.4}
A particularly important case where $Y$ is separated over $S$ is the case where $Y=S[T]=S\otimes_{\bb{Z}}\Spec(\bb{Z}[T])$ ($T$ an indeterminate).
There is then a bijective correspondence between the pseudo-$S$-morphisms from $X$ to $Y$ and the \emph{pseudo-functions} on $X$ \sref{IV.20.2.8}, by virtue of the definition of a product of $\bb{Z}$-preschemes.
The pseudo-morphisms from $X$ to $Y$ relative to $S$ are then identified, by virtue of \sref{IV.20.5.3}, with the \emph{pseudo-functions $\varphi$ on $X$ such that $\operatorname{dom}(\varphi)$ is universally schematically dense relative to $S$}.
The sheaf $\mathscr{P}\!s.\!\operatorname{hom}_{X/S}(X,Y)$ is a \emph{sub-sheaf of rings} of $\sh{M}'_X$, which we denote by $\sh{M}'_{X/S}$.

We then set $\operatorname{Ps.hom}_{X/S}(X,Y)=M'(X/S)$ and say that its elements are the \emph{pseudo-functions on $X$ relative to $S$}.
\end{env}

\begin{env}[20.5.5]
\label{IV.20.5.5}
Let $X$, $Y$, $Z$ be three $S$-preschemes and $f:Y\to Z$ an $S$-morphism.
In the argument of \sref{IV.20.3.1}, we can replace ``schematically dense'' everywhere by ``universally schematically dense relative to $S$''; for every pseudo-morphism $\omega\in\operatorname{Ps.hom}_{X/S}(X,Y)$, the morphisms $f\circ u$, where $u$ runs through $\omega$, therefore belong to one and the same equivalence class (for the relation defined in \sref{IV.20.5.1}), which we
\oldpage[IV-4]{251}
call the \emph{composite} of $f$ and $\omega$, a pseudo-morphism from $X$ to $Z$ relative to $S$, and which we denote by $f\circ\omega$.
If $g:Z\to T$ is a morphism, it is immediate that $g\circ(f\circ\omega)=(g\circ f)\circ\omega$.
\end{env}

\begin{env}[20.5.6]
\label{IV.20.5.6}
Suppose $Y$ \emph{separated} over $S$, and let $\omega$ be a pseudo-morphism from $X$ to $Y$ relative to $S$.
If $f:X'\to X$ is an $S$-morphism such that $f^{-1}(\operatorname{dom}_S(\omega))$ is \emph{universally schematically dense in $X'$ relative to $S$}, it follows from \sref{IV.20.3.2} that the pseudo-$S$-morphism $\omega\circ f$ has a domain that is universally schematically dense relative to $S$, and hence \sref{IV.20.5.3} it can be regarded as a pseudo-morphism \emph{relative to $S$}.
When $X'$ is flat and locally of finite presentation over $S$, the condition that $f^{-1}(\operatorname{dom}_S(\omega))$ be universally schematically dense relative to $S$ is again equivalent to saying that, for every $s\in S$, $(f^{-1}(\operatorname{dom}_S(\omega)))_s$ (notation of \sref{IV.11.10.10}) is schematically dense in $X'_s$, or again, denoting by $f_s:X'_s\to X_s$ the morphism deduced from $f$ by base change, that the inverse image under $f_s$ of $(\operatorname{dom}_S(\omega))_s$ is schematically dense in $X'_s$.
This last condition will in particular be fulfilled if, for every $s\in S$, $X_s$, $X'_s$ and $f_s$ satisfy one of the three conditions~(i),~(ii),~(iii) of \sref{IV.20.3.5}.
\end{env}

\begin{env}[20.5.7]
\label{IV.20.5.7}
Suppose now that $X$ and $X'$ are both $S$-preschemes \emph{flat and locally of finite presentation over $S$}, and that $f:X'\to X$ is a \emph{flat} $S$-morphism (or, what amounts to the same \sref{IV.11.3.10}, that for every $s\in S$, $f_s:X'_s\to X_s$ is a flat morphism).
Then, for every open subset $U$ of $X$ and every open subset $V\subset U$ universally schematically dense in $U$ relative to $S$, it follows from \sref{IV.11.10.5} and \sref{IV.11.10.9} that $f^{-1}(V)$ is universally schematically dense in $f^{-1}(U)$ relative to $S$.
For every pseudo-morphism $\omega$ from $X$ to $Y$ relative to $S$, it follows from \sref{IV.20.3.4} that the pseudo-$S$-morphism $\omega\circ f$ is defined and is a pseudo-morphism from $X'$ to $Y$ relative to $S$, \emph{even when $Y$ is not assumed separated over $S$}.
We deduce that in this case, for every $S$-morphism $g:Y\to Z$, $(g\circ\omega)\circ f$ is again defined and equal to $g\circ(\omega\circ f)$ (with the definitions of \sref{IV.20.5.1}), and is therefore a pseudo-morphism \emph{relative to $S$}.
\end{env}

\begin{env}[20.5.8]
\label{IV.20.5.8}
Let $X$ be an $S$-prescheme, $S'\to S$ an arbitrary morphism, $X'=X_{(S')}$, and $g:X'\to X$ the canonical projection.
Then, for every open subset $U$ of $X$ and every open subset $V\subset U$ universally schematically dense in $U$ relative to $S$, $V'=g^{-1}(V)$ is universally schematically dense in $U'=g^{-1}(U)$ relative to $S'$, by definition \sref{IV.11.10.8}.
Now let $\omega$ be a pseudo-morphism from $X$ to an $S$-prescheme $Y$ relative to $S$; if $u_1$, $u_2$ are $S$-morphisms of the class $\omega$, defined respectively on open subsets $U_1$, $U_2$ of $X$ universally schematically dense in $X$ relative to $S$, it follows from the above that the morphisms $u'_1=u_1\circ(g|g^{-1}(U_1))$ and $u'_2=u_2\circ(g|g^{-1}(U_2))$ coincide on an open subset $U'_3$ universally schematically dense relative to $S'$.
Now, if $Y'=Y_{(S')}$ and if $h:Y'\to Y$ is the canonical projection, $u'_1$ and $u'_2$ factor canonically as $u'_1=h\circ v_1$, $u'_2=h\circ v_2$, and $v_1$ and $v_2$ are two $S'$-morphisms into $Y'$ which coincide on $U'_3$.
We therefore see that, as $u_1$ runs through the class $\omega$, the corresponding $S'$-morphisms $v_1$ belong to one and the same pseudo-morphism relative to $S'$, called the \emph{inverse image} of $\omega$ under the base-change morphism $S'\to S$ and denoted by $\omega_{(S')}$.
It is clear that if $S''\to S'$ is a second morphism, we have $(\omega_{(S')})_{(S'')}=\omega_{(S'')}$ (for the composite base-change morphism $S''\to S'\to S$).
\end{env}

\oldpage[IV-4]{252}

\subsection{Relative meromorphic functions}
\EGAReaderTarget{ega4.c04.subsec.20-6}{IV.20.6}

\begin{env}[20.6.1]\label{IV.20.6.1}
\EGAReaderTarget{ega4.c04.statement.20-6-1}{IV.20.6.1}
Let $S$ be a prescheme and $X$ an $S$-prescheme which is \emph{flat and locally of finite presentation} over $S$; for every $s\in S$, we will denote by $X_s$ the fibre at the point $s$ of the structural morphism $X\to S$.
In general, if $\varphi$ is a meromorphic function on $X$, it is not possible to associate with it in a ``natural'' way, for each $s\in S$, an ``induced'' meromorphic function on $X_s$ \sref{IV.20.1.11}.
For every open subset $U$ of $X$, denote by $T_{X/S}(U)$ the set of sections $t\in\Gamma(U,\sh{O}_X)$ such that, for every $s\in S$, the image $t_s$ of $t$ under the canonical homomorphism $\Gamma(U,\sh{O}_X)\to\Gamma(U\cap X_s,\sh{O}_{X_s})$ is a \emph{regular} section; this implies moreover, by the equivalence of a) and b) in \sref{IV.11.3.7}, that $t$ is itself a \emph{regular} section.
It is clear that $U\mapsto T_{X/S}(U)$ is a \emph{sub-sheaf} of the sheaf of sets $\sh{S}=\sh{S}(\sh{O}_X)$ (notation of \sref{IV.20.1.3}), which we denote by $\sh{T}_{X/S}$; we set
\[
\label{IV.20.6.1.1}
  \sh{M}_{X/S}=\sh{O}_X[\sh{T}_{X/S}^{-1}]
  \tag{20.6.1.1}
\]
(notation of \sref{IV.20.1.1}) and we say that this sheaf of rings is the \emph{sheaf of germs of meromorphic functions on $X$ relative to $S$}; its sections over $X$ are called the \emph{meromorphic functions on $X$ relative to $S$}, and their set is denoted by $M(X/S)$.
It is clear that $\sh{M}_{X/S}$ is a \emph{sub-sheaf} of $\sh{M}_X=\sh{O}_X[\sh{S}^{-1}]$; for every meromorphic function $\varphi\in M(X/S)$ and every $s\in S$, the inverse image of $\varphi$ under the canonical injection morphism $j_s:X_s\to X$ is then defined \sref{IV.20.1.11}, and is denoted by $\varphi_s$.
\end{env}

\begin{env}[20.6.2]\label{IV.20.6.2}
\EGAReaderTarget{ega4.c04.statement.20-6-2}{IV.20.6.2}
Now let $\sh{F}$ be a quasi-coherent $\sh{O}_X$-module; we set
\[
\label{IV.20.6.2.1}
  \sh{M}_{X/S}(\sh{F})=\sh{F}[\sh{T}_{X/S}^{-1}]=\sh{F}\otimes_{\sh{O}_X}\sh{M}_{X/S};
  \tag{20.6.2.1}
\]
the sections of $\sh{M}_{X/S}(\sh{F})$ are called \emph{meromorphic sections of $\sh{F}$ over $X$, relative to $S$}, and their set is denoted by $M(X/S,F)$.
The canonical homomorphism $\sh{F}\to\sh{M}_{X/S}(\sh{F})$ is not necessarily injective; when it is, we say that $\sh{F}$ is \emph{torsion-free relative to $S$}: this means that, for every open subset $U$ of $X$ and every section $t\in T_{X/S}(U)$, $t$ is $(\sh{F}|U)$-regular; this condition is fulfilled \emph{a fortiori} when $\sh{F}$ is \emph{strictly torsion-free} \sref{IV.20.1.5}.
In the latter case, it follows immediately from the definitions \sref{IV.20.1.2} that the canonical homomorphism of $\sh{O}_X$-modules
\[
\label{IV.20.6.2.2}
  \sh{M}_{X/S}(\sh{F})\to\sh{M}_X(\sh{F})
  \tag{20.6.2.2}
\]
is \emph{injective}, so that the meromorphic sections of $\sh{F}$ relative to $S$ are meromorphic sections of $\sh{F}$ in the sense of \sref{IV.20.1.3}.
\end{env}

\begin{proposition}[20.6.3]
\label{IV.20.6.3}
The image under the injective homomorphism \sref{IV.20.2.10.1} of the sub-sheaf $\sh{M}_{X/S}$ of $\sh{M}_X$ is the sub-sheaf $\sh{M}'_{X/S}$ of pseudo-functions on $X$ relative to $S$ (i.e.\ of the pseudo-functions whose domain of definition is universally schematically dense relative to $S$ \sref{IV.20.5.4}).
\end{proposition}

\begin{proof}
We may evidently restrict to proving that the image of $M(X/S)$ under the canonical homomorphism $M(X)\to M'(X)$ is equal to $M'(X/S)$; the proposition is then a consequence of the following more general proposition:
\end{proof}

\oldpage[IV-4]{253}
\begin{proposition}[20.6.4]
\label{IV.20.6.4}
Let $\sh{F}$ be a quasi-coherent $\sh{O}_X$-module of finite presentation which is strictly torsion-free.
Then, for a meromorphic section $\varphi$ of $\sh{F}$ over $X$ to be a meromorphic section relative to $S$, it is necessary and sufficient that $\operatorname{dom}(\varphi)$ be universally schematically dense relative to $S$.
\end{proposition}

\begin{proof}
The necessity of the condition follows from \sref{IV.20.2.15} applied to each canonical injection $X_s\to X$ ($s\in S$), taking into account \sref{IV.11.10.9}.
To see that the condition is sufficient, we must prove that for every $x\in X$ there exist an open neighbourhood $U$ of $x$ in $X$ and a section of $\sh{M}_{X/S}(\sh{F})$ over $U$ whose restriction to $U\cap\operatorname{dom}(\varphi)$ coincides with $\varphi$ on a schematically dense open subset of $U\cap\operatorname{dom}(\varphi)$.
Consider the ideal of denominators $\sh{J}$ of $\varphi$ \sref{IV.20.2.14}, which is quasi-coherent and defines a closed sub-prescheme of $X$ whose underlying space is $X-\operatorname{dom}(\varphi)$.
By hypothesis, if $s$ is the image of $x$ in $S$, $\operatorname{dom}(\varphi)\cap X_s$ is schematically dense in the locally Noetherian prescheme $X_s$, and therefore \sref{IV.20.2.13}[(iv)] contains $\operatorname{Ass}(\sh{O}_{X_s})$; this therefore implies that the ideal $\sh{J}_x$ of $\sh{O}_{X,x}$ has an image in $\sh{O}_{X_s,x}=\sh{O}_{X,x}\otimes_{\sh{O}_{S,s}}k(s)$ which is contained in none of the prime ideals $\mathfrak{p}_i\in\operatorname{Ass}(\sh{O}_{X_s,x})$ (which are finite in number); hence (Bourbaki, \emph{Alg.~comm.}, chap.~II, \textsection1, n\textsuperscript{o}~1, prop.~2) there exists an element $t_x\in\sh{J}_x$ whose image in $\sh{O}_{X_s,x}$ belongs to none of the $\mathfrak{p}_i$, and is consequently regular in this Noetherian ring.
Let $t$ be a section of $\sh{J}$ over an affine open neighbourhood $U$ of $x$ whose germ at the point $x$ is $t_x$; since $X$ is flat and locally of finite presentation over $S$, we may suppose \sref{IV.11.3.8} that $t$ is a \emph{regular} section of $\sh{O}_X$ over $U$ and that, for every $s'\in S$, the image of $t$ in $\Gamma(U\cap X_{s'},\sh{O}_{X_{s'}})$ is also \emph{regular}; in other words, we have $t\in T_{X/S}(U)$.
But then, by definition of $\sh{J}$, since $\sh{F}$ is strictly torsion-free, $t(\varphi|(U\cap\operatorname{dom}(\varphi)))$ is a section $u$ of $\sh{F}$ over $U\cap\operatorname{dom}(\varphi)$; on the other hand, $U\cap\operatorname{dom}(\varphi)$ contains the open set $U_t$ of points $x'\in U$ where $t(x')\neq0$, and the latter contains $x$ and is schematically dense in $U$ \sref{IV.20.2.9}.
We therefore see that on $U_t$, $\varphi$ coincides with the restriction to $U_t$ of the section $u/t$ of $\sh{M}_{X/S}(\sh{F})$ over $U$.
Q.E.D.
\end{proof}

\begin{remarks}[20.6.5]\label{IV.20.6.5}
\EGAReaderTarget{ega4.c04.statement.20-6-5}{IV.20.6.5}
\begin{enumerate}[label=(\roman*)]
  \item Let $\varphi$ be a meromorphic function on $X$ relative to $S$, so that for every $s\in S$, $\varphi_s$ is a meromorphic function on $X_s$ \sref{IV.20.6.1}; by virtue of \sref{IV.20.1.11.1}, we have
  \[
  \label{IV.20.6.5.1}
    \operatorname{dom}(\varphi)\cap X_s\subset\operatorname{dom}(\varphi_s).
    \tag{20.6.5.1}
  \]
  But it should be noted that, even when $S$ is the spectrum of a discrete valuation ring $A$ and $X=S[T]$ ($T$ an indeterminate), the two sides of \sref{IV.20.6.5.1} are not necessarily equal: for example, if $\pi$ is a uniformizer of $A$, it is immediate that $\varphi=\pi/T$ is a meromorphic function on $S$ relative to $S$, because, if $a$ and $b$ are the closed point and the generic point of $S$, $T$ is regular in $\Gamma(X_a,\sh{O}_{X_a})=k[T]$ and in $\Gamma(X_b,\sh{O}_{X_b})=K[T]$, $k$ and $K$ being the residue field and the field of fractions of $A$.
  We have $\operatorname{dom}(\varphi)=\mathrm{D}(T)$ in $X$, but $\operatorname{dom}(\varphi_a)=X_a$ since $\varphi_a=0$.

  \item For a meromorphic function $\varphi$ relative to $S$ to be \emph{invertible in the ring} $M(X/S)$, it is necessary and sufficient that, for every $s\in S$, $\varphi_s$ be invertible in $M(X_s)$ (in other words, that $\varphi_s$ be a regular meromorphic function on $X_s$ \sref{IV.20.1.8}).
  The condition is
  \oldpage[IV-4]{254}
  indeed trivially necessary.
  Conversely, if it is satisfied, and if $x$ is an arbitrary point of $X$ and $s$ its image in $S$, there exist by hypothesis an open neighbourhood $U$ of $x$ in $X$ and two sections $u$, $t$ of $\sh{O}_X$ over $U$ such that $t\in T_{X/S}(U)$ and $\varphi|U=u/t$; the hypothesis implies that, if $u_s$ is the image of $u$ in $\Gamma(U\cap X_s,\sh{O}_{X_s})$, then $u_s$ is regular at the point $x$.
  By shrinking $U$, we may therefore suppose, by virtue of \sref{IV.11.3.8}, that $u\in T_{X/S}(U)$, whence the conclusion.

  When $\varphi$ is invertible in $M(X/S)$, we also say that $\varphi$ is a \emph{regular meromorphic function relative to $S$}.
  Note that $\varphi\in M(X/S)$ may be invertible in $M(X)$ (in other words, be \emph{regular} in the sense of \sref{IV.20.1.8}) without being invertible in $M(X/S)$, as the example of~(i) immediately shows.

  \item Let $\sh{L}$ be an \emph{invertible} $\sh{O}_X$-module, and let $\varphi$ be a \emph{regular} meromorphic section of $\sh{L}$ over $X$ \sref{IV.20.1.8}; we say that $\varphi$ is \emph{regular relative to $S$} if, for every open subset $U$ of $X$ such that $\sh{L}|U$ is isomorphic to $\sh{O}_U$, $\varphi|U$ corresponds to an element of $\Gamma(U,\sh{M}_X)$ which is \emph{regular relative to $S$}; by virtue of~(ii), it is immediate that for this it is necessary and sufficient that, for every $s\in S$, $\varphi_s$ be a \emph{regular} meromorphic section \sref{IV.20.1.8} of the invertible $\sh{O}_{X_s}$-module $\sh{L}_s=\sh{L}\otimes_{\sh{O}_X}k(s)$.
  If $\varphi'$ is the inverse of $\varphi$ in $\sh{L}^{-1}$ \sref{IV.20.1.10}, then $\varphi'$ is also regular relative to $S$.
  If $\sh{L}_1$ is a second invertible $\sh{O}_X$-module and $\varphi_1$ a meromorphic section of $\sh{L}_1$ over $X$ which is regular relative to $S$, then $\varphi\otimes\varphi_1$ is a meromorphic section of $\sh{L}\otimes\sh{L}_1$ over $X$ which is regular relative to $S$.
\end{enumerate}
\end{remarks}

\begin{proposition}[20.6.6]
\label{IV.20.6.6}
Let $X$ be an $S$-prescheme flat and locally of finite presentation over $S$, and $\sh{F}$ a locally free $\sh{O}_X$-module of finite type; for every $s\in S$, denote by $X_s$ the fibre at the point $s$ of the structural morphism $f:X\to S$.
Let $\varphi$ be a meromorphic section of $\sh{F}$ over $X$, relative to $S$, and suppose that $\varphi$ is defined at all the points $x\in X$ such that $\operatorname{prof}(\sh{O}_{X_{f(x)},x})=1$.
Then $\varphi$ is everywhere defined.
\end{proposition}

\begin{proof}
By hypothesis, $\operatorname{dom}(\varphi)\cap X_s$ is schematically dense in $X_s$ for every $s\in S$, so it contains the points $x$ of $X_s$ such that $\operatorname{prof}(\sh{O}_{X_s,x})=0$ \sref{IV.5.10.2}; the hypothesis therefore means that, if we set $Z=X-\operatorname{dom}(\varphi)$, we have $\operatorname{prof}(\sh{O}_{X_{f(x)},x})\geq2$ at all the points of $Z$.
It therefore suffices to apply \sref{IV.19.9.8}.
\end{proof}

\begin{env}[20.6.7]\label{IV.20.6.7}
\EGAReaderTarget{ega4.c04.statement.20-6-7}{IV.20.6.7}
Let $X$, $X'$ be two $S$-preschemes flat and locally of finite presentation over $S$, and $f=(\psi,\theta):X'\to X$ an $S$-morphism.
For every open subset $U$ of $X$, denote by $T_f(U)$ the set of sections $t\in T_{X/S}(U)$ whose image in $\Gamma(f^{-1}(U),\sh{O}_{X'})$ belongs to $T_{X'/S}(f^{-1}(U))$; it is immediate that $U\mapsto T_f(U)$ is a sub-sheaf of the sheaf of sets $\sh{T}_{X/S}$, which we denote by $\sh{T}_f$.
We set $\sh{M}_{X/S,f}=\sh{O}_X[\sh{T}_f^{-1}]$; this is a sub-sheaf of rings of $\sh{M}_{X/S}$, and from $\theta^\#: \psi^*(\sh{O}_X)\to\sh{O}_{X'}$ we canonically deduce a homomorphism of sheaves of rings $\theta'^\#: \psi^*(\sh{M}_{X/S,f})\to\sh{M}_{X'/S}$ extending $\theta^\#$.
If a meromorphic function $\varphi$ on $X$ relative to $S$ is \emph{a section of $\sh{M}_{X/S,f}$}, then $\Gamma(\theta'^\#)(\varphi)$ is a meromorphic function on $X'$, called the \emph{inverse image of $\varphi$ under $f$} and denoted by $\varphi\circ f$ if no confusion can result.
The definitions of \sref{IV.20.1.11} relating to $\sh{O}_X$-modules are immediately extended in the same way.
\end{env}

\begin{proposition}[20.6.8]
\label{IV.20.6.8}
With the notation of \sref{IV.20.6.7}, if the $S$-morphism $f:X'\to X$ is \emph{flat}, we have $\sh{M}_{X/S,f}=\sh{M}_{X/S}$, and the homomorphism $\varphi\mapsto\varphi\circ f$ is defined on the whole of $M(X/S)$.
\end{proposition}

\begin{proof}
\oldpage[IV-4]{255}
Indeed, the hypothesis implies, by virtue of \sref{IV.11.3.10}, that for every $s\in S$, $f_s:X'_s\to X_s$ is flat; hence, for every section $t\in T_{X/S}(U)$, if $t'$ is its inverse image in $\Gamma(f^{-1}(U),\sh{O}_{X'})$, then $t'_{s}$, which is the inverse image of $t_s$, is a \emph{regular} section of $\sh{O}_{X'_s}$ over $f^{-1}(U)\cap X'_s$, by virtue of \sref{IV.20.1.12}; we conclude that, by definition, $t'\in T_{X'/S}(f^{-1}(U))$, whence the proposition.
\end{proof}

From this we deduce, as in \sref{IV.20.1.12}, a canonical homomorphism of $\sh{O}_{X'}$-algebras
\[
\label{IV.20.6.8.1}
  f^*(\sh{M}_{X/S})\to\sh{M}_{X'/S}.
  \tag{20.6.8.1}
\]

\begin{env}[20.6.9]
\label{IV.20.6.9}
Finally, consider an \emph{arbitrary} morphism $S'\to S$, and set $X'=X\times_S S'$, which is flat and locally of finite presentation over $S'$; let $p:X'\to X$ be the canonical projection.
Let $U$ be an open subset of $X$, $t$ a section belonging to $T_{X/S}(U)$, and $t'$ its image in $\Gamma(p^{-1}(U),\sh{O}_{X'})$; for every $s'\in S'$, if $s\in S$ is the image of $s'$, we have $X'_{s'}=X_s\otimes_{k(s)}k(s')$, so the morphism $X'_{s'}\to X_s$ is flat, and consequently \sref{IV.20.1.12} the inverse image $t'_{s'}$ of $t_s$ in $\Gamma(p^{-1}(U)\cap X'_{s'},\sh{O}_{X'_{s'}})$ is regular; this proves that we have $t'\in T_{X'/S'}(p^{-1}(U))$.
This allows us to define canonically, as in \sref{IV.20.6.8}, a canonical homomorphism of $\sh{O}_{X'}$-algebras
\[
\label{IV.20.6.9.1}
  p^*(\sh{M}_{X/S})\to\sh{M}_{X'/S'}.
  \tag{20.6.9.1}
\]

With the identification permitted by \sref{IV.20.6.3}, this notion of base change for relative meromorphic functions is a particular case of the analogous notion for relative pseudo-morphisms \sref{IV.20.5.8}.
\end{env}
